% Generated by src/make_supplementary.py from saved results; do not edit by hand.
\documentclass[10pt]{article}
\usepackage[utf8]{inputenc}
\usepackage[T1]{fontenc}
\usepackage{lmodern}
\usepackage[a4paper,margin=2cm]{geometry}
\usepackage{booktabs}
\usepackage{longtable}
\usepackage{array}
\usepackage{amsmath}
\usepackage[hidelinks]{hyperref}
\makeatletter\g@addto@macro\UrlBreaks{\do\_\do\-}\makeatother
\renewcommand{\thetable}{S\arabic{table}}
\renewcommand{\thesection}{S\arabic{section}}
\setcounter{secnumdepth}{2}
\title{Supplementary Information\\[4pt] \large Rendering gain and sequence modeling in 12-lead ECG-image classification: a controlled comparison with external validation}
\author{Miriam Istanboli and Mohamed Elgendi}
\date{}
\begin{document}
\maketitle
\noindent Every number in this document was generated from the saved run records and analysis outputs by
\nolinkurl{src/make_supplementary.py}. PTB-XL test-fold (fold 10) results are exploratory, because that fold had been inspected
before the models were specified; the external cohorts are the confirmatory test. A completed TRIPOD+AI checklist
\cite{collins2024tripodai} is provided as a separate file.
\tableofcontents

\section{Cohorts and labels}

\begingroup\footnotesize
\setlength{\tabcolsep}{3pt}
\begin{longtable}{lrrrrrrr}
\caption{PTB-XL cohort by official fold: records, patients and positives per superclass.}\label{tab:s_cohort}\\
\toprule
Partition & Records & Patients & NORM & MI & STTC & CD & HYP \\
\midrule
\endfirsthead
\toprule
Partition & Records & Patients & NORM & MI & STTC & CD & HYP \\
\midrule
\endhead
\bottomrule
\endfoot
Training (folds 1--8) & 17{,}084 & 14{,}823 & 7{,}596 & 4{,}379 & 4{,}186 & 3{,}907 & 2{,}119 \\
Validation (fold 9) & 2{,}146 & 1{,}917 & 955 & 540 & 528 & 495 & 268 \\
Test (fold 10) & 2{,}158 & 1{,}877 & 963 & 550 & 521 & 496 & 262 \\
\midrule
Total & 21{,}388 & 18{,}617 & 9{,}514 & 5{,}469 & 5{,}235 & 4{,}898 & 2{,}649 \\
\end{longtable}
\endgroup
{\footnotesize PTB-XL v1.0.3 contains 21{,}799 records from 18{,}869 patients; 411 records without any diagnostic superclass were excluded, and every included record was rendered and used. Labels overlap, so positives do not sum to the number of records.\par}
\medskip

\begingroup\footnotesize
\setlength{\tabcolsep}{3pt}
\begin{longtable}{lrr}
\caption{External cohorts: inclusion flow and positives per superclass.}\label{tab:s_flow}\\
\toprule
Step & Georgia, USA & China \\
\midrule
\endfirsthead
\toprule
Step & Georgia, USA & China \\
\midrule
\endhead
\bottomrule
\endfoot
Records in the Challenge 2021 training set & 10{,}344 & 45{,}152 \\
Excluded: signal file failed PhysioNet SHA-256 (amendment) & 0 & 2 \\
Excluded: not 12 standard leads & 0 & 0 \\
Excluded: not 500 Hz & 0 & 0 \\
Excluded: not in millivolts & 0 & 0 \\
Excluded: shorter than 10 s & 52 & 0 \\
Excluded: no mapped superclass & 4{,}245 & 36{,}146 \\
Included & 6{,}047 & 9{,}004 \\
\midrule
Positives, NORM & 0$^*$ & 0$^*$ \\
Positives, MI & 7$^*$ & 169 \\
Positives, STTC & 4{,}059 & 5{,}670 \\
Positives, CD & 1{,}823 & 3{,}857 \\
Positives, HYP & 2{,}088 & 770 \\
\end{longtable}
\endgroup
{\footnotesize $^*$Not evaluated in that cohort: a label needs at least 50 positives and 50 negatives (prespecified). Steps are applied in the order shown; each record is counted at its first failed step. China: Chapman-Shaoxing and Ningbo databases. Source: experiments/external/$<$cohort$>$/inventory.json.\par}
\medskip

\begingroup\scriptsize
\setlength{\tabcolsep}{3pt}
\begin{longtable}{lp{4.6cm}rrrrrrl}
\caption{SNOMED CT code, clinical name and PTB-XL superclass mapping used for the external cohorts.}\label{tab:s_mapping}\\
\toprule
SNOMED CT & Clinical name (abbreviation) & PTB-XL records & NORM & MI & STTC & CD & HYP & Assigned \\
\midrule
\endfirsthead
\toprule
SNOMED CT & Clinical name (abbreviation) & PTB-XL records & NORM & MI & STTC & CD & HYP & Assigned \\
\midrule
\endhead
\bottomrule
\endfoot
426783006 & sinus rhythm (NSR) & 18{,}058 & 0.53 & 0.22 & 0.19 & 0.20 & 0.11 & -- \\
164865005 & myocardial infarction (MI) & 5{,}248 & 0.00 & \textbf{1.00} & 0.24 & 0.33 & 0.15 & MI \\
39732003 & left axis deviation (LAD) & 5{,}146 & 0.20 & 0.44 & 0.23 & 0.45 & 0.12 & -- \\
164951009 & abnormal QRS (abQRS) & 3{,}389 & 0.08 & 0.81 & 0.21 & 0.28 & 0.14 & -- \\
164873001 & left ventricular hypertrophy (LVH) & 2{,}354 & 0.05 & 0.30 & 0.62 & 0.21 & \textbf{0.91} & HYP \\
164934002 & t wave abnormal (TAb) & 2{,}341 & 0.05 & 0.16 & 0.88 & 0.16 & 0.12 & -- \\
164861001 & myocardial ischemia (MIs) & 2{,}171 & 0.00 & 0.38 & \textbf{1.00} & 0.25 & 0.58 & STTC \\
445118002 & left anterior fascicular block (LAnFB) & 1{,}624 & 0.01 & 0.46 & 0.24 & \textbf{1.00} & 0.13 & CD \\
164889003 & atrial fibrillation (AF) & 1{,}514 & 0.02 & 0.40 & 0.60 & 0.30 & 0.18 & -- \\
164884008 & ventricular ectopics (VEB) & 1{,}153 & 0.17 & 0.42 & 0.38 & 0.33 & 0.20 & -- \\
713426002 & incomplete right bundle branch block (IRBBB) & 1{,}118 & 0.22 & 0.24 & 0.15 & \textbf{1.00} & 0.13 & CD \\
429622005 & st depression (STD) & 1{,}009 & 0.06 & 0.38 & 0.74 & 0.23 & 0.41 & -- \\
427084000 & sinus tachycardia (STach) & 826 & 0.29 & 0.36 & 0.32 & 0.29 & 0.20 & -- \\
270492004 & 1st degree av block (IAVB) & 795 & 0.05 & 0.47 & 0.26 & \textbf{1.00} & 0.22 & CD \\
698252002 & nonspecific intraventricular conduction disorder (NSIVCB) & 787 & 0.13 & 0.47 & 0.36 & \textbf{1.00} & 0.23 & CD \\
427393009 & sinus arrhythmia (SA) & 772 & 0.55 & 0.21 & 0.18 & 0.19 & 0.10 & -- \\
55930002 & s t changes (STC) & 768 & 0.04 & 0.14 & \textbf{1.00} & 0.18 & 0.13 & STTC \\
426177001 & sinus bradycardia (SB) & 637 & 0.55 & 0.17 & 0.18 & 0.21 & 0.11 & -- \\
164917005 & qwave abnormal (QAb) & 548 & 0.10 & 0.80 & 0.31 & 0.35 & 0.18 & -- \\
713427006 & complete right bundle branch block (CRBBB) & 542 & 0.00 & 0.53 & 0.05 & \textbf{1.00} & 0.15 & CD \\
164909002 & left bundle branch block (LBBB) & 536 & 0.00 & 0.07 & 0.04 & \textbf{1.00} & 0.14 & CD \\
67741000119109 & left atrial enlargement (LAE) & 426 & 0.01 & 0.38 & 0.36 & 0.58 & \textbf{1.00} & HYP \\
284470004 & premature atrial contraction (PAC) & 398 & 0.32 & 0.29 & 0.35 & 0.32 & 0.21 & -- \\
428750005 & nonspecific st t abnormality (NSSTTA) & 381 & 0.00 & 0.41 & \textbf{1.00} & 0.23 & 0.21 & STTC \\
54329005 & anterior myocardial infarction (AnMI) & 353 & 0.00 & \textbf{1.00} & 0.27 & 0.34 & 0.15 & MI \\
47665007 & right axis deviation (RAD) & 343 & 0.19 & 0.47 & 0.16 & 0.59 & 0.16 & -- \\
164947007 & prolonged pr interval (LPR) & 340 & 0.24 & 0.39 & 0.34 & 0.55 & 0.26 & -- \\
10370003 & pacing rhythm (PR) & 294 & 0.00 & 0.01 & 0.01 & 0.01 & 0.01 & -- \\
59931005 & t wave inversion (TInv) & 294 & 0.02 & 0.63 & 0.66 & 0.27 & 0.35 & -- \\
425419005 & inferior ischaemia (IIs) & 218 & 0.00 & 0.19 & \textbf{1.00} & 0.18 & 0.13 & STTC \\
251146004 & low qrs voltages (LQRSV) & 182 & 0.43 & 0.25 & 0.27 & 0.15 & 0.03 & -- \\
445211001 & left posterior fascicular block (LPFB) & 177 & 0.04 & 0.47 & 0.12 & \textbf{1.00} & 0.23 & CD \\
63593006 & supraventricular premature beats (SVPB) & 157 & 0.23 & 0.38 & 0.26 & 0.36 & 0.15 & -- \\
251200008 & indeterminate cardiac axis (ICA) & 155 & 0.37 & 0.18 & 0.23 & 0.34 & 0.18 & -- \\
425623009 & lateral ischaemia (LIs) & 140 & 0.00 & 0.64 & \textbf{1.00} & 0.28 & 0.08 & STTC \\
89792004 & right ventricular hypertrophy (RVH) & 126 & 0.00 & 0.17 & 0.06 & \textbf{0.94} & \textbf{1.00} & CD, HYP \\
111975006 & prolonged qt interval (LQT) & 118 & 0.00 & 0.19 & \textbf{0.99} & 0.19 & 0.10 & STTC \\
446358003 & right atrial hypertrophy (RAH) & 99 & 0.01 & 0.29 & 0.26 & 0.45 & \textbf{1.00} & HYP \\
11157007 & ventricular bigeminy (VBig) & 82 & 0.16 & 0.33 & 0.27 & 0.26 & 0.16 & -- \\
74390002 & wolff parkinson white pattern (WPW) & 79 & 0.00 & 0.08 & 0.00 & \textbf{1.00} & 0.00 & CD \\
251120003 & incomplete left bundle branch block (ILBBB) & 77 & 0.00 & 0.52 & 0.51 & \textbf{1.00} & 0.38 & CD \\
164890007 & atrial flutter (AFL) & 73 & 0.00 & 0.12 & 0.21 & 0.18 & 0.04 & -- \\
426434006 & anterior ischemia (AnMIs) & 44 & 0.00 & 0.64 & \textbf{1.00} & 0.32 & 0.09 & STTC \\
266249003 & ventricular hypertrophy (VH) & 29 & 0.00 & 0.17 & 0.48 & 0.21 & \textbf{1.00} & HYP \\
164931005 & st elevation (STE) & 28 & 0.46 & 0.11 & 0.36 & 0.18 & 0.07 & -- \\
426761007 & supraventricular tachycardia (SVT) & 27 & 0.04 & 0.15 & 0.67 & 0.30 & 0.15 & -- \\
67198005 & paroxysmal supraventricular tachycardia (PSVT) & 24 & 0.00 & 0.00 & 0.33 & 0.04 & 0.00 & -- \\
251180001 & ventricular trigeminy (VTrig) & 20 & 0.30 & 0.30 & 0.30 & 0.15 & 0.15 & -- \\
27885002 & complete heart block (CHB) & 16 & 0.00 & 0.56 & 0.25 & \textbf{1.00} & 0.19 & CD \\
195042002 & 2nd degree av block (IIAVB) & 14 & 0.07 & 0.07 & 0.50 & \textbf{1.00} & 0.21 & CD \\
\end{longtable}
\endgroup
{\footnotesize Rule fixed before the mapping was computed: share >= 0.9 among codes on >= 10 PTB-XL records. Shares are the fraction of PTB-XL records carrying the SNOMED CT code (Challenge organizers' encoding of PTB-XL) that have each superclass; bold shares meet the 90\% rule. 25 of 50 codes were assigned. Headers checked: 21{,}799; age/sex mismatches: 0. Clinical names (added 15 November 2026): the diagnosis name and abbreviation that the PhysioNet/CinC Challenge 2021 organizers give for each code (files dx\_mapping\_scored.csv and dx\_mapping\_unscored.csv of the repository physionetchallenges/evaluation-2021), lower-cased as in the source. The name was compared with the SNOMED CT concept name from an independent terminology server for 50 of the 50 codes (49 agree in meaning, 1 differ; the rest could not be checked); sources and notes per code are in experiments/external/snomed\_code\_names.csv. The mapping itself is data-driven and was not reviewed by a cardiologist (Section~\ref{sec:s_review}).\par}
\medskip

\section{Every run: results and selected epochs}

\begingroup\scriptsize
\setlength{\tabcolsep}{3pt}
\begin{longtable}{lllllrrrr}
\caption{Every PTB-XL run: hardware, selected epoch, validation and test (fold 10, exploratory) results.}\label{tab:s_runs}\\
\toprule
Pages & Model & Seed & Device & Selected epoch & Val. AUROC & Test AUROC & Test AP & Test F1 \\
\midrule
\endfirsthead
\toprule
Pages & Model & Seed & Device & Selected epoch & Val. AUROC & Test AUROC & Test AP & Test F1 \\
\midrule
\endhead
\bottomrule
\endfoot
Adaptive & CNN-only & 42 & CPU & 19 of 20 & 0.845 & 0.834 & 0.631 & 0.496 \\
Adaptive & CNN-only & 43 & GPU & 20 of 20 & 0.843 & 0.827 & 0.617 & 0.494 \\
Adaptive & CNN-only & 44 & GPU & 19 of 20 & 0.835 & 0.826 & 0.620 & 0.481 \\
Adaptive & CNN-only & 45 & GPU & 18 of 20 & 0.841 & 0.832 & 0.628 & 0.491 \\
Adaptive & CNN-only & 46 & GPU & 20 of 20 & 0.840 & 0.832 & 0.628 & 0.487 \\
Adaptive & CNN-only, capacity-matched & 42 & GPU & 20 of 20 & 0.874 & 0.865 & 0.696 & 0.567 \\
Adaptive & CNN-only, capacity-matched & 43 & GPU & 20 of 20 & 0.880 & 0.868 & 0.705 & 0.604 \\
Adaptive & CNN-only, capacity-matched & 44 & GPU & 18 of 20 & 0.874 & 0.862 & 0.689 & 0.578 \\
Adaptive & CNN-only, capacity-matched & 45 & GPU & 18 of 20 & 0.863 & 0.859 & 0.686 & 0.571 \\
Adaptive & CNN-only, capacity-matched & 46 & GPU & 19 of 20 & 0.882 & 0.872 & 0.711 & 0.604 \\
Adaptive & CNN-LSTM & 42 & CPU & 15 of 19 & 0.882 & 0.870 & 0.704 & 0.613 \\
Adaptive & CNN-LSTM & 43 & GPU & 16 of 20 & 0.894 & 0.881 & 0.734 & 0.646 \\
Adaptive & CNN-LSTM & 44 & GPU & 18 of 20 & 0.902 & 0.886 & 0.740 & 0.653 \\
Adaptive & CNN-LSTM & 45 & GPU & 18 of 20 & 0.888 & 0.872 & 0.719 & 0.636 \\
Adaptive & CNN-LSTM & 46 & GPU & 16 of 20 & 0.890 & 0.873 & 0.713 & 0.638 \\
Adaptive & CNN + temporal conv. & 42 & CPU & 16 of 20 & 0.889 & 0.880 & 0.732 & 0.641 \\
Adaptive & CNN + temporal conv. & 43 & GPU & 19 of 20 & 0.885 & 0.879 & 0.729 & 0.632 \\
Adaptive & CNN + temporal conv. & 44 & GPU & 20 of 20 & 0.888 & 0.879 & 0.730 & 0.639 \\
Adaptive & CNN + temporal conv. & 45 & GPU & 16 of 20 & 0.893 & 0.883 & 0.733 & 0.625 \\
Adaptive & CNN + temporal conv. & 46 & GPU & 20 of 20 & 0.885 & 0.875 & 0.723 & 0.617 \\
Adaptive & CNN + small ViT & 42 & CPU & 19 of 20 & 0.902 & 0.890 & 0.752 & 0.672 \\
Adaptive & CNN + small ViT & 43 & GPU & 18 of 20 & 0.895 & 0.888 & 0.748 & 0.674 \\
Adaptive & CNN + small ViT & 44 & GPU & 13 of 17 & 0.898 & 0.886 & 0.746 & 0.667 \\
Adaptive & CNN + small ViT & 45 & GPU & 13 of 17 & 0.892 & 0.876 & 0.720 & 0.647 \\
Adaptive & CNN + small ViT & 46 & GPU & 17 of 20 & 0.893 & 0.887 & 0.741 & 0.667 \\
\midrule
Fixed & CNN-only & 42 & GPU & 19 of 20 & 0.868 & 0.854 & 0.680 & 0.567 \\
Fixed & CNN-only & 43 & GPU & 20 of 20 & 0.876 & 0.867 & 0.706 & 0.596 \\
Fixed & CNN-only & 44 & GPU & 20 of 20 & 0.866 & 0.861 & 0.696 & 0.582 \\
Fixed & CNN-only & 45 & GPU & 18 of 20 & 0.870 & 0.856 & 0.680 & 0.577 \\
Fixed & CNN-only & 46 & GPU & 18 of 20 & 0.875 & 0.863 & 0.697 & 0.588 \\
Fixed & CNN-only, capacity-matched$^\P$ & 42 & GPU & 18 of 20 & 0.892 & 0.882 & 0.741 & 0.643 \\
Fixed & CNN-only, capacity-matched$^\P$ & 43 & GPU & 18 of 20 & 0.888 & 0.875 & 0.725 & 0.628 \\
Fixed & CNN-only, capacity-matched$^\P$ & 44 & GPU & 17 of 20 & 0.889 & 0.880 & 0.733 & 0.640 \\
Fixed & CNN-only, capacity-matched$^\P$ & 45 & GPU & 20 of 20 & 0.897 & 0.886 & 0.744 & 0.645 \\
Fixed & CNN-only, capacity-matched$^\P$ & 46 & GPU & 20 of 20 & 0.899 & 0.890 & 0.754 & 0.653 \\
Fixed & CNN-LSTM & 42 & GPU & 16 of 20 & 0.898 & 0.884 & 0.739 & 0.641 \\
Fixed & CNN-LSTM & 43 & GPU & 18 of 20 & 0.900 & 0.886 & 0.744 & 0.664 \\
Fixed & CNN-LSTM & 44 & GPU & 20 of 20 & 0.899 & 0.887 & 0.745 & 0.668 \\
Fixed & CNN-LSTM & 45 & GPU & 20 of 20 & 0.906 & 0.896 & 0.763 & 0.675 \\
Fixed & CNN-LSTM & 46 & GPU & 16 of 20 & 0.898 & 0.892 & 0.754 & 0.674 \\
Fixed & CNN + temporal conv. & 42 & GPU & 18 of 20 & 0.912 & 0.896 & 0.766 & 0.679 \\
Fixed & CNN + temporal conv. & 43 & GPU & 18 of 20 & 0.901 & 0.894 & 0.763 & 0.678 \\
Fixed & CNN + temporal conv. & 44 & GPU & 18 of 20 & 0.906 & 0.894 & 0.763 & 0.678 \\
Fixed & CNN + temporal conv. & 45 & GPU & 20 of 20 & 0.904 & 0.895 & 0.763 & 0.674 \\
Fixed & CNN + temporal conv. & 46 & GPU & 16 of 20 & 0.909 & 0.896 & 0.760 & 0.682 \\
Fixed & CNN + small ViT & 42 & GPU & 15 of 19 & 0.903 & 0.899 & 0.766 & 0.684 \\
Fixed & CNN + small ViT & 43 & GPU & 18 of 20 & 0.899 & 0.882 & 0.737 & 0.669 \\
Fixed & CNN + small ViT & 44 & GPU & 13 of 17 & 0.906 & 0.894 & 0.752 & 0.677 \\
Fixed & CNN + small ViT & 45 & GPU & 14 of 18 & 0.905 & 0.894 & 0.749 & 0.684 \\
Fixed & CNN + small ViT & 46 & GPU & 16 of 20 & 0.899 & 0.887 & 0.748 & 0.677 \\
\midrule
Adaptive & PCA + LDA (grid $\le$256) & -- & CPU & -- & 0.801 & 0.788 & 0.570 & 0.403 \\
Adaptive & PCA + LDA (grid 256--1024) & -- & CPU & -- & 0.806 & 0.794 & 0.576 & 0.427 \\
Adaptive & PCA + linear SVM & -- & CPU & -- & 0.806 & 0.793 & 0.577 & 0.411 \\
\midrule
Fixed & PCA + LDA (grid 256--1024) & -- & CPU & -- & 0.790 & 0.781 & 0.576 & 0.450 \\
\midrule
Signal & 1D ResNet (raw signals) & 42 & GPU & 15 of 19 & 0.932 & 0.929 & 0.823 & 0.731 \\
Signal & 1D ResNet (raw signals) & 43 & GPU & 12 of 16 & 0.931 & 0.928 & 0.825 & 0.743 \\
Signal & 1D ResNet (raw signals) & 44 & GPU & 12 of 16 & 0.931 & 0.928 & 0.819 & 0.729 \\
Signal & 1D ResNet (raw signals) & 45 & GPU & 12 of 16 & 0.931 & 0.925 & 0.819 & 0.736 \\
Signal & 1D ResNet (raw signals) & 46 & GPU & 10 of 14 & 0.932 & 0.927 & 0.822 & 0.727 \\
\end{longtable}
\endgroup
{\footnotesize Macro averages over the five labels; F1 at the fixed threshold 0.5. Selected epoch: epoch of the checkpoint chosen on validation, of the epochs run (maximum 20; fewer means early stopping). CPU: laptop CPU (Mac, seed 42 of the four original adaptive-gain networks); GPU: NVIDIA RTX 5070 Ti. PCA references are single fits. Seed 42 of CNN-LSTM, temporal convolution and ViT with adaptive gain are the original v1 runs. $^\P$Additional analysis (amendment of 15 November 2026, Section~\ref{sec:s_review}): the capacity-matched head trained on fixed-gain pages, added after all other results were known. \par}
\medskip

\begingroup\scriptsize
\setlength{\tabcolsep}{3pt}
\begin{longtable}{lllrrrrr}
\caption{Every PTB-XL run: test AUROC per superclass.}\label{tab:s_runs_labels}\\
\toprule
Pages & Model & Seed & NORM & MI & STTC & CD & HYP \\
\midrule
\endfirsthead
\toprule
Pages & Model & Seed & NORM & MI & STTC & CD & HYP \\
\midrule
\endhead
\bottomrule
\endfoot
Adaptive & CNN-only & 42 & 0.894 & 0.801 & 0.859 & 0.815 & 0.802 \\
Adaptive & CNN-only & 43 & 0.883 & 0.788 & 0.864 & 0.809 & 0.791 \\
Adaptive & CNN-only & 44 & 0.884 & 0.790 & 0.867 & 0.817 & 0.769 \\
Adaptive & CNN-only & 45 & 0.889 & 0.799 & 0.867 & 0.812 & 0.793 \\
Adaptive & CNN-only & 46 & 0.884 & 0.796 & 0.883 & 0.804 & 0.792 \\
Adaptive & CNN-only, capacity-matched & 42 & 0.905 & 0.854 & 0.896 & 0.859 & 0.811 \\
Adaptive & CNN-only, capacity-matched & 43 & 0.908 & 0.839 & 0.911 & 0.840 & 0.841 \\
Adaptive & CNN-only, capacity-matched & 44 & 0.900 & 0.810 & 0.900 & 0.842 & 0.860 \\
Adaptive & CNN-only, capacity-matched & 45 & 0.910 & 0.828 & 0.914 & 0.838 & 0.805 \\
Adaptive & CNN-only, capacity-matched & 46 & 0.908 & 0.836 & 0.907 & 0.856 & 0.851 \\
Adaptive & CNN-LSTM & 42 & 0.913 & 0.829 & 0.903 & 0.846 & 0.858 \\
Adaptive & CNN-LSTM & 43 & 0.919 & 0.840 & 0.910 & 0.866 & 0.868 \\
Adaptive & CNN-LSTM & 44 & 0.923 & 0.860 & 0.908 & 0.866 & 0.874 \\
Adaptive & CNN-LSTM & 45 & 0.917 & 0.845 & 0.899 & 0.845 & 0.857 \\
Adaptive & CNN-LSTM & 46 & 0.915 & 0.819 & 0.912 & 0.846 & 0.874 \\
Adaptive & CNN + temporal conv. & 42 & 0.920 & 0.867 & 0.912 & 0.856 & 0.846 \\
Adaptive & CNN + temporal conv. & 43 & 0.927 & 0.859 & 0.915 & 0.864 & 0.830 \\
Adaptive & CNN + temporal conv. & 44 & 0.921 & 0.845 & 0.915 & 0.863 & 0.851 \\
Adaptive & CNN + temporal conv. & 45 & 0.923 & 0.867 & 0.916 & 0.866 & 0.841 \\
Adaptive & CNN + temporal conv. & 46 & 0.920 & 0.854 & 0.914 & 0.851 & 0.839 \\
Adaptive & CNN + small ViT & 42 & 0.923 & 0.869 & 0.916 & 0.872 & 0.871 \\
Adaptive & CNN + small ViT & 43 & 0.920 & 0.862 & 0.916 & 0.874 & 0.870 \\
Adaptive & CNN + small ViT & 44 & 0.919 & 0.859 & 0.914 & 0.866 & 0.870 \\
Adaptive & CNN + small ViT & 45 & 0.912 & 0.831 & 0.906 & 0.874 & 0.860 \\
Adaptive & CNN + small ViT & 46 & 0.921 & 0.866 & 0.911 & 0.867 & 0.872 \\
\midrule
Fixed & CNN-only & 42 & 0.904 & 0.805 & 0.887 & 0.817 & 0.854 \\
Fixed & CNN-only & 43 & 0.908 & 0.824 & 0.905 & 0.836 & 0.864 \\
Fixed & CNN-only & 44 & 0.902 & 0.810 & 0.891 & 0.836 & 0.864 \\
Fixed & CNN-only & 45 & 0.905 & 0.808 & 0.884 & 0.820 & 0.864 \\
Fixed & CNN-only & 46 & 0.912 & 0.812 & 0.890 & 0.835 & 0.868 \\
Fixed & CNN-only, capacity-matched$^\P$ & 42 & 0.916 & 0.849 & 0.908 & 0.853 & 0.886 \\
Fixed & CNN-only, capacity-matched$^\P$ & 43 & 0.911 & 0.838 & 0.905 & 0.847 & 0.872 \\
Fixed & CNN-only, capacity-matched$^\P$ & 44 & 0.916 & 0.843 & 0.911 & 0.856 & 0.874 \\
Fixed & CNN-only, capacity-matched$^\P$ & 45 & 0.917 & 0.872 & 0.909 & 0.861 & 0.872 \\
Fixed & CNN-only, capacity-matched$^\P$ & 46 & 0.925 & 0.863 & 0.919 & 0.863 & 0.878 \\
Fixed & CNN-LSTM & 42 & 0.923 & 0.845 & 0.914 & 0.854 & 0.883 \\
Fixed & CNN-LSTM & 43 & 0.923 & 0.845 & 0.914 & 0.871 & 0.875 \\
Fixed & CNN-LSTM & 44 & 0.920 & 0.852 & 0.921 & 0.869 & 0.873 \\
Fixed & CNN-LSTM & 45 & 0.929 & 0.874 & 0.922 & 0.870 & 0.882 \\
Fixed & CNN-LSTM & 46 & 0.924 & 0.856 & 0.920 & 0.872 & 0.887 \\
Fixed & CNN + temporal conv. & 42 & 0.928 & 0.876 & 0.922 & 0.871 & 0.882 \\
Fixed & CNN + temporal conv. & 43 & 0.930 & 0.872 & 0.921 & 0.867 & 0.880 \\
Fixed & CNN + temporal conv. & 44 & 0.927 & 0.865 & 0.916 & 0.879 & 0.886 \\
Fixed & CNN + temporal conv. & 45 & 0.929 & 0.876 & 0.924 & 0.863 & 0.883 \\
Fixed & CNN + temporal conv. & 46 & 0.927 & 0.866 & 0.920 & 0.877 & 0.889 \\
Fixed & CNN + small ViT & 42 & 0.930 & 0.863 & 0.922 & 0.889 & 0.888 \\
Fixed & CNN + small ViT & 43 & 0.917 & 0.847 & 0.915 & 0.861 & 0.867 \\
Fixed & CNN + small ViT & 44 & 0.926 & 0.868 & 0.917 & 0.877 & 0.881 \\
Fixed & CNN + small ViT & 45 & 0.927 & 0.878 & 0.919 & 0.864 & 0.883 \\
Fixed & CNN + small ViT & 46 & 0.918 & 0.857 & 0.920 & 0.871 & 0.872 \\
\midrule
Adaptive & PCA + LDA (grid $\le$256) & -- & 0.813 & 0.771 & 0.743 & 0.799 & 0.816 \\
Adaptive & PCA + LDA (grid 256--1024) & -- & 0.823 & 0.772 & 0.754 & 0.801 & 0.820 \\
Adaptive & PCA + linear SVM & -- & 0.822 & 0.770 & 0.753 & 0.801 & 0.818 \\
\midrule
Fixed & PCA + LDA (grid 256--1024) & -- & 0.798 & 0.752 & 0.720 & 0.799 & 0.837 \\
\midrule
Signal & 1D ResNet (raw signals) & 42 & 0.946 & 0.934 & 0.935 & 0.925 & 0.905 \\
Signal & 1D ResNet (raw signals) & 43 & 0.946 & 0.932 & 0.936 & 0.921 & 0.905 \\
Signal & 1D ResNet (raw signals) & 44 & 0.946 & 0.934 & 0.934 & 0.921 & 0.903 \\
Signal & 1D ResNet (raw signals) & 45 & 0.945 & 0.932 & 0.935 & 0.919 & 0.895 \\
Signal & 1D ResNet (raw signals) & 46 & 0.945 & 0.935 & 0.934 & 0.924 & 0.899 \\
\end{longtable}
\endgroup
{\footnotesize $^\P$Additional analysis (amendment of 15 November 2026, Section~\ref{sec:s_review}): the capacity-matched head trained on fixed-gain pages, added after all other results were known. \par}
\medskip

\begingroup\footnotesize
\setlength{\tabcolsep}{3pt}
\begin{longtable}{llccc}
\caption{Selected (validation-best) epoch per run, by model.}\label{tab:s_epochs}\\
\toprule
Pages & Model & Selected epoch per seed & Selected $\ge$18 & Stopped early \\
\midrule
\endfirsthead
\toprule
Pages & Model & Selected epoch per seed & Selected $\ge$18 & Stopped early \\
\midrule
\endhead
\bottomrule
\endfoot
Adaptive & CNN-only & 19, 20, 19, 18, 20 & 5 of 5 & 0 of 5 \\
Adaptive & CNN-only, capacity-matched & 20, 20, 18, 18, 19 & 5 of 5 & 0 of 5 \\
Adaptive & CNN-LSTM & 15, 16, 18, 18, 16 & 2 of 5 & 1 of 5 \\
Adaptive & CNN + temporal conv. & 16, 19, 20, 16, 20 & 3 of 5 & 0 of 5 \\
Adaptive & CNN + small ViT & 19, 18, 13, 13, 17 & 2 of 5 & 2 of 5 \\
Fixed & CNN-only & 19, 20, 20, 18, 18 & 5 of 5 & 0 of 5 \\
Fixed & CNN-only, capacity-matched$^\P$ & 18, 18, 17, 20, 20 & 4 of 5 & 0 of 5 \\
Fixed & CNN-LSTM & 16, 18, 20, 20, 16 & 3 of 5 & 0 of 5 \\
Fixed & CNN + temporal conv. & 18, 18, 18, 20, 16 & 4 of 5 & 0 of 5 \\
Fixed & CNN + small ViT & 15, 18, 13, 14, 16 & 1 of 5 & 3 of 5 \\
Signal & 1D ResNet (raw signals) & 15, 12, 12, 12, 10 & 0 of 5 & 5 of 5 \\
\end{longtable}
\endgroup
{\footnotesize Seeds 42--46 in order. Image networks selecting epoch 18--20 of 20: 17 of 25 adaptive-gain runs and 17 of 25 fixed-gain runs. The epoch budget was prespecified and not tuned; a cosine schedule that ends at epoch 20 favours late selection, but late selection may also mean that some models were still improving. $^\P$Additional analysis (amendment of 15 November 2026, Section~\ref{sec:s_review}): the capacity-matched head trained on fixed-gain pages, added after all other results were known. The fixed-gain count includes its five runs.\par}
\medskip

\section{Sensitivity analyses with GPU-trained seeds only}

\begingroup\scriptsize
\setlength{\tabcolsep}{3pt}
\begin{longtable}{lcccc}
\caption{PTB-XL head comparisons (adaptive gain) with all seeds and with GPU-trained seeds only: difference in mean-over-seeds macro AUROC [95\% CI] and Holm-adjusted $p$.}\label{tab:s_gpu_pairs}\\
\toprule
Comparison & \multicolumn{2}{c}{Seeds 42--46} & \multicolumn{2}{c}{Seeds 43--46, GPU only} \\
\midrule
\endfirsthead
\toprule
Comparison & \multicolumn{2}{c}{Seeds 42--46} & \multicolumn{2}{c}{Seeds 43--46, GPU only} \\
\midrule
\endhead
\bottomrule
\endfoot
\multicolumn{5}{l}{\emph{Test fold (exploratory)}} \\
CNN-only, capacity-matched $-$ CNN-only & +0.035 [+0.028, +0.042] & $\le$0.002 & +0.036 [+0.028, +0.044] & $\le$0.002 \\
CNN-LSTM $-$ CNN-only & +0.046 [+0.036, +0.057] & $\le$0.002 & +0.049 [+0.039, +0.060] & $\le$0.002 \\
CNN + temporal conv $-$ CNN-only & +0.049 [+0.042, +0.056] & $\le$0.002 & +0.050 [+0.042, +0.057] & $\le$0.002 \\
CNN + small ViT $-$ CNN-only & +0.055 [+0.047, +0.064] & $\le$0.002 & +0.055 [+0.045, +0.065] & $\le$0.002 \\
CNN-LSTM $-$ CNN-only, capacity-matched & +0.011 [+0.003, +0.020] & 0.022 & +0.013 [+0.004, +0.022] & 0.021 \\
CNN + temporal conv $-$ CNN-only, capacity-matched & +0.014 [+0.007, +0.021] & 0.004 & +0.014 [+0.006, +0.022] & 0.009 \\
CNN + small ViT $-$ CNN-only, capacity-matched & +0.020 [+0.014, +0.026] & $\le$0.002 & +0.019 [+0.013, +0.025] & $\le$0.002 \\
CNN + temporal conv $-$ CNN-LSTM & +0.003 [$-$0.005, +0.010] & 0.452 & +0.001 [$-$0.007, +0.009] & 0.846 \\
CNN + small ViT $-$ CNN-LSTM & +0.009 [+0.002, +0.017] & 0.044 & +0.006 [$-$0.000, +0.013] & 0.200 \\
CNN + small ViT $-$ CNN + temporal conv & +0.006 [$-$0.001, +0.013] & 0.186 & +0.005 [$-$0.003, +0.013] & 0.408 \\
\midrule
\multicolumn{5}{l}{\emph{Validation fold}} \\
CNN-only, capacity-matched $-$ CNN-only & +0.034 [+0.025, +0.042] & $\le$0.002 & +0.035 [+0.025, +0.044] & $\le$0.002 \\
CNN-LSTM $-$ CNN-only & +0.050 [+0.040, +0.061] & $\le$0.002 & +0.054 [+0.045, +0.064] & $\le$0.002 \\
CNN + temporal conv $-$ CNN-only & +0.047 [+0.041, +0.054] & $\le$0.002 & +0.048 [+0.041, +0.055] & $\le$0.002 \\
CNN + small ViT $-$ CNN-only & +0.055 [+0.047, +0.063] & $\le$0.002 & +0.054 [+0.047, +0.063] & $\le$0.002 \\
CNN-LSTM $-$ CNN-only, capacity-matched & +0.017 [+0.009, +0.025] & $\le$0.002 & +0.019 [+0.010, +0.028] & $\le$0.002 \\
CNN + temporal conv $-$ CNN-only, capacity-matched & +0.014 [+0.005, +0.023] & 0.007 & +0.013 [+0.003, +0.025] & 0.022 \\
CNN + small ViT $-$ CNN-only, capacity-matched & +0.021 [+0.013, +0.029] & $\le$0.002 & +0.020 [+0.011, +0.029] & $\le$0.002 \\
CNN + temporal conv $-$ CNN-LSTM & $-$0.003 [$-$0.011, +0.005] & 0.499 & $-$0.006 [$-$0.013, +0.002] & 0.307 \\
CNN + small ViT $-$ CNN-LSTM & +0.005 [$-$0.002, +0.013] & 0.499 & +0.001 [$-$0.004, +0.006] & 0.759 \\
CNN + small ViT $-$ CNN + temporal conv & +0.008 [+0.002, +0.014] & 0.034 & +0.007 [+0.000, +0.013] & 0.147 \\
\end{longtable}
\endgroup
{\footnotesize Hierarchical bootstrap (seeds, then patients; 10{,}000 resamples), Holm over the 10 pairs within each fold and seed set. $p$ printed as ``$\le$'' is at the resolution limit of the bootstrap: no replicate crossed zero, and no smaller value is claimed (Section~\ref{sec:s_review}). Seeds 43--46 exclude the CPU-trained seed 42 of the original networks.\par}
\medskip

\begingroup\footnotesize
\setlength{\tabcolsep}{3pt}
\begin{longtable}{llcccc}
\caption{PTB-XL test: fixed-gain minus adaptive-gain training with all seeds and with GPU-trained seeds only [95\% CI], Holm-adjusted $p$.}\label{tab:s_gpu_fixed}\\
\toprule
Architecture & AUROC & \multicolumn{2}{c}{Seeds 42--46} & \multicolumn{2}{c}{Seeds 43--46, GPU only} \\
\midrule
\endfirsthead
\toprule
Architecture & AUROC & \multicolumn{2}{c}{Seeds 42--46} & \multicolumn{2}{c}{Seeds 43--46, GPU only} \\
\midrule
\endhead
\bottomrule
\endfoot
CNN-only & macro & +0.030 [+0.021, +0.039] & $\le$0.002 & +0.033 [+0.024, +0.041] & $\le$0.002 \\
 & HYP & +0.073 [+0.050, +0.098] & $\le$0.002 & +0.079 [+0.057, +0.103] & $\le$0.002 \\
CNN + temporal conv. & macro & +0.016 [+0.011, +0.021] & $\le$0.002 & +0.016 [+0.010, +0.022] & $\le$0.002 \\
 & HYP & +0.043 [+0.028, +0.058] & $\le$0.002 & +0.044 [+0.029, +0.060] & $\le$0.002 \\
CNN-LSTM & macro & +0.012 [+0.004, +0.021] & 0.017 & +0.012 [+0.002, +0.022] & 0.084 \\
 & HYP & +0.014 [+0.001, +0.028] & 0.106 & +0.011 [$-$0.002, +0.026] & 0.310 \\
CNN + small ViT & macro & +0.006 [$-$0.003, +0.014] & 0.251 & +0.005 [$-$0.005, +0.015] & 0.514 \\
 & HYP & +0.010 [$-$0.003, +0.023] & 0.251 & +0.008 [$-$0.006, +0.023] & 0.514 \\
\end{longtable}
\endgroup
{\footnotesize Hierarchical bootstrap (seeds, then patients; 10{,}000 resamples); Holm over the 8 comparisons within each seed set. $p$ printed as ``$\le$'' is at the resolution limit of the bootstrap: no replicate crossed zero, and no smaller value is claimed (Section~\ref{sec:s_review}). With seeds 42--46, the seed-42 pair combines a CPU-trained adaptive model with a GPU-trained fixed-gain model; the GPU-only set removes that hardware difference.\par}
\medskip

\section{Calibration and validation-chosen thresholds}

\begingroup\scriptsize
\setlength{\tabcolsep}{3pt}
\begin{longtable}{lllccccc}
\caption{Calibration and validation-chosen thresholds (mean $\pm$ SD over runs).}\label{tab:s_calibration}\\
\toprule
Pages & Model & Runs & Val. ECE & Test ECE & Test Brier & Test F1 (0.5) & Test F1 (val.\ thr.) \\
\midrule
\endfirsthead
\toprule
Pages & Model & Runs & Val. ECE & Test ECE & Test Brier & Test F1 (0.5) & Test F1 (val.\ thr.) \\
\midrule
\endhead
\bottomrule
\endfoot
Adaptive & CNN-only & 5 & 0.024 $\pm$ 0.003 & 0.026 $\pm$ 0.002 & 0.125 $\pm$ 0.001 & 0.490 $\pm$ 0.006 & 0.602 $\pm$ 0.006 \\
Adaptive & CNN-only, capacity-matched & 5 & 0.024 $\pm$ 0.002 & 0.028 $\pm$ 0.003 & 0.112 $\pm$ 0.002 & 0.585 $\pm$ 0.018 & 0.647 $\pm$ 0.007 \\
Adaptive & CNN-LSTM & 5 & 0.020 $\pm$ 0.002 & 0.026 $\pm$ 0.003 & 0.107 $\pm$ 0.003 & 0.637 $\pm$ 0.015 & 0.673 $\pm$ 0.013 \\
Adaptive & CNN + temporal conv. & 5 & 0.029 $\pm$ 0.006 & 0.033 $\pm$ 0.006 & 0.104 $\pm$ 0.001 & 0.631 $\pm$ 0.010 & 0.672 $\pm$ 0.006 \\
Adaptive & CNN + small ViT & 5 & 0.041 $\pm$ 0.005 & 0.046 $\pm$ 0.003 & 0.105 $\pm$ 0.002 & 0.666 $\pm$ 0.011 & 0.684 $\pm$ 0.005 \\
\midrule
Fixed & CNN-only & 5 & 0.022 $\pm$ 0.002 & 0.022 $\pm$ 0.004 & 0.114 $\pm$ 0.002 & 0.582 $\pm$ 0.011 & 0.645 $\pm$ 0.009 \\
Fixed & CNN-only, capacity-matched$^\P$ & 5 & 0.021 $\pm$ 0.002 & 0.022 $\pm$ 0.003 & 0.104 $\pm$ 0.003 & 0.642 $\pm$ 0.009 & 0.679 $\pm$ 0.008 \\
Fixed & CNN-LSTM & 5 & 0.022 $\pm$ 0.002 & 0.026 $\pm$ 0.001 & 0.101 $\pm$ 0.002 & 0.664 $\pm$ 0.014 & 0.690 $\pm$ 0.012 \\
Fixed & CNN + temporal conv. & 5 & 0.025 $\pm$ 0.002 & 0.028 $\pm$ 0.003 & 0.099 $\pm$ 0.001 & 0.678 $\pm$ 0.003 & 0.702 $\pm$ 0.005 \\
Fixed & CNN + small ViT & 5 & 0.040 $\pm$ 0.007 & 0.045 $\pm$ 0.006 & 0.103 $\pm$ 0.003 & 0.678 $\pm$ 0.006 & 0.695 $\pm$ 0.012 \\
\midrule
Adaptive & PCA + LDA (grid $\le$256) & 1 & 0.026 & 0.022 & 0.141 & 0.403 & 0.564 \\
Adaptive & PCA + LDA (grid 256--1024) & 1 & 0.032 & 0.026 & 0.140 & 0.427 & 0.570 \\
Adaptive & PCA + linear SVM & 1 & 0.029 & 0.022 & 0.140 & 0.411 & 0.570 \\
\midrule
Fixed & PCA + LDA (grid 256--1024) & 1 & 0.034 & 0.033 & 0.143 & 0.450 & 0.565 \\
\midrule
Signal & 1D ResNet (raw signals) & 5 & 0.026 $\pm$ 0.005 & 0.028 $\pm$ 0.005 & 0.082 $\pm$ 0.001 & 0.733 $\pm$ 0.006 & 0.749 $\pm$ 0.003 \\
\end{longtable}
\endgroup
{\footnotesize Macro averages over the five labels. ECE: expected calibration error with 10 equal-width bins. Validation thresholds: per label, the threshold in 0.01--0.99 that maximizes validation F1, applied unchanged to test. The fixed 0.5 rule remains primary.\par}
\medskip

\begingroup\scriptsize
\setlength{\tabcolsep}{3pt}
\begin{longtable}{lllrrrrrrr}
\caption{Validation-chosen thresholds per label and their effect on test (mean over runs).}\label{tab:s_thresholds}\\
\toprule
Pages & Model & Label & Thr. & Sens.\ 0.5 & Spec.\ 0.5 & Sens.\ thr. & Spec.\ thr. & F1 0.5 & F1 thr. \\
\midrule
\endfirsthead
\toprule
Pages & Model & Label & Thr. & Sens.\ 0.5 & Spec.\ 0.5 & Sens.\ thr. & Spec.\ thr. & F1 0.5 & F1 thr. \\
\midrule
\endhead
\bottomrule
\endfoot
Adaptive & CNN-only & NORM & 0.42 & 0.83 & 0.78 & 0.87 & 0.74 & 0.79 & 0.79 \\
 &  & MI & 0.31 & 0.34 & 0.93 & 0.63 & 0.80 & 0.44 & 0.57 \\
 &  & STTC & 0.32 & 0.47 & 0.94 & 0.69 & 0.85 & 0.57 & 0.64 \\
 &  & CD & 0.32 & 0.41 & 0.96 & 0.58 & 0.87 & 0.53 & 0.58 \\
 &  & HYP & 0.20 & 0.07 & 0.99 & 0.51 & 0.88 & 0.13 & 0.43 \\
\midrule
Adaptive & CNN-only, capacity-matched & NORM & 0.43 & 0.87 & 0.78 & 0.89 & 0.74 & 0.81 & 0.81 \\
 &  & MI & 0.29 & 0.44 & 0.93 & 0.69 & 0.81 & 0.53 & 0.61 \\
 &  & STTC & 0.35 & 0.56 & 0.94 & 0.73 & 0.89 & 0.64 & 0.70 \\
 &  & CD & 0.35 & 0.47 & 0.96 & 0.58 & 0.92 & 0.59 & 0.63 \\
 &  & HYP & 0.25 & 0.24 & 0.98 & 0.54 & 0.91 & 0.35 & 0.49 \\
\midrule
Adaptive & CNN-LSTM & NORM & 0.41 & 0.87 & 0.81 & 0.90 & 0.77 & 0.82 & 0.83 \\
 &  & MI & 0.31 & 0.48 & 0.93 & 0.67 & 0.83 & 0.56 & 0.62 \\
 &  & STTC & 0.35 & 0.61 & 0.94 & 0.73 & 0.89 & 0.68 & 0.70 \\
 &  & CD & 0.32 & 0.51 & 0.96 & 0.62 & 0.91 & 0.62 & 0.64 \\
 &  & HYP & 0.30 & 0.39 & 0.98 & 0.57 & 0.94 & 0.50 & 0.57 \\
\midrule
Adaptive & CNN + temporal conv. & NORM & 0.50 & 0.90 & 0.79 & 0.90 & 0.79 & 0.83 & 0.83 \\
 &  & MI & 0.30 & 0.50 & 0.94 & 0.68 & 0.85 & 0.59 & 0.64 \\
 &  & STTC & 0.35 & 0.64 & 0.94 & 0.74 & 0.90 & 0.70 & 0.72 \\
 &  & CD & 0.32 & 0.53 & 0.96 & 0.61 & 0.92 & 0.64 & 0.65 \\
 &  & HYP & 0.25 & 0.27 & 0.98 & 0.53 & 0.92 & 0.39 & 0.51 \\
\midrule
Adaptive & CNN + small ViT & NORM & 0.46 & 0.88 & 0.80 & 0.89 & 0.78 & 0.83 & 0.83 \\
 &  & MI & 0.28 & 0.49 & 0.94 & 0.67 & 0.86 & 0.59 & 0.64 \\
 &  & STTC & 0.41 & 0.71 & 0.92 & 0.76 & 0.89 & 0.72 & 0.72 \\
 &  & CD & 0.33 & 0.57 & 0.95 & 0.65 & 0.92 & 0.66 & 0.68 \\
 &  & HYP & 0.29 & 0.48 & 0.96 & 0.60 & 0.92 & 0.53 & 0.55 \\
\midrule
Fixed & CNN-only & NORM & 0.44 & 0.84 & 0.81 & 0.87 & 0.77 & 0.81 & 0.81 \\
 &  & MI & 0.27 & 0.37 & 0.94 & 0.69 & 0.78 & 0.48 & 0.59 \\
 &  & STTC & 0.34 & 0.55 & 0.93 & 0.74 & 0.85 & 0.62 & 0.67 \\
 &  & CD & 0.30 & 0.46 & 0.96 & 0.63 & 0.88 & 0.58 & 0.62 \\
 &  & HYP & 0.24 & 0.31 & 0.98 & 0.60 & 0.91 & 0.43 & 0.54 \\
\midrule
Fixed & CNN-only, capacity-matched$^\P$ & NORM & 0.47 & 0.86 & 0.81 & 0.87 & 0.80 & 0.82 & 0.82 \\
 &  & MI & 0.27 & 0.46 & 0.94 & 0.73 & 0.80 & 0.56 & 0.63 \\
 &  & STTC & 0.33 & 0.61 & 0.94 & 0.79 & 0.87 & 0.68 & 0.72 \\
 &  & CD & 0.33 & 0.52 & 0.96 & 0.64 & 0.90 & 0.63 & 0.65 \\
 &  & HYP & 0.26 & 0.41 & 0.98 & 0.62 & 0.92 & 0.52 & 0.57 \\
\midrule
Fixed & CNN-LSTM & NORM & 0.42 & 0.87 & 0.82 & 0.90 & 0.79 & 0.83 & 0.83 \\
 &  & MI & 0.32 & 0.49 & 0.93 & 0.69 & 0.85 & 0.58 & 0.64 \\
 &  & STTC & 0.32 & 0.67 & 0.94 & 0.79 & 0.88 & 0.72 & 0.74 \\
 &  & CD & 0.38 & 0.56 & 0.95 & 0.63 & 0.93 & 0.65 & 0.67 \\
 &  & HYP & 0.28 & 0.44 & 0.97 & 0.60 & 0.93 & 0.54 & 0.57 \\
\midrule
Fixed & CNN + temporal conv. & NORM & 0.46 & 0.87 & 0.83 & 0.88 & 0.82 & 0.83 & 0.84 \\
 &  & MI & 0.32 & 0.55 & 0.93 & 0.71 & 0.85 & 0.62 & 0.66 \\
 &  & STTC & 0.33 & 0.67 & 0.94 & 0.78 & 0.90 & 0.72 & 0.74 \\
 &  & CD & 0.35 & 0.58 & 0.95 & 0.65 & 0.92 & 0.67 & 0.68 \\
 &  & HYP & 0.29 & 0.46 & 0.97 & 0.61 & 0.94 & 0.54 & 0.59 \\
\midrule
Fixed & CNN + small ViT & NORM & 0.49 & 0.89 & 0.81 & 0.89 & 0.80 & 0.83 & 0.83 \\
 &  & MI & 0.30 & 0.53 & 0.92 & 0.69 & 0.85 & 0.60 & 0.65 \\
 &  & STTC & 0.36 & 0.69 & 0.93 & 0.77 & 0.89 & 0.72 & 0.74 \\
 &  & CD & 0.28 & 0.56 & 0.96 & 0.67 & 0.91 & 0.66 & 0.68 \\
 &  & HYP & 0.33 & 0.53 & 0.96 & 0.61 & 0.93 & 0.58 & 0.58 \\
\midrule
Adaptive & PCA + LDA (grid $\le$256) & NORM & 0.42 & 0.71 & 0.76 & 0.81 & 0.66 & 0.71 & 0.73 \\
 &  & MI & 0.29 & 0.27 & 0.95 & 0.61 & 0.79 & 0.38 & 0.55 \\
 &  & STTC & 0.26 & 0.17 & 0.96 & 0.62 & 0.74 & 0.26 & 0.51 \\
 &  & CD & 0.26 & 0.29 & 0.97 & 0.62 & 0.84 & 0.42 & 0.57 \\
 &  & HYP & 0.24 & 0.15 & 0.99 & 0.53 & 0.89 & 0.25 & 0.46 \\
\midrule
Adaptive & PCA + LDA (grid 256--1024) & NORM & 0.41 & 0.72 & 0.77 & 0.83 & 0.66 & 0.72 & 0.73 \\
 &  & MI & 0.25 & 0.26 & 0.95 & 0.69 & 0.73 & 0.37 & 0.56 \\
 &  & STTC & 0.26 & 0.19 & 0.96 & 0.61 & 0.76 & 0.29 & 0.52 \\
 &  & CD & 0.28 & 0.33 & 0.97 & 0.58 & 0.87 & 0.46 & 0.58 \\
 &  & HYP & 0.27 & 0.19 & 0.99 & 0.49 & 0.91 & 0.30 & 0.46 \\
\midrule
Adaptive & PCA + linear SVM & NORM & 0.40 & 0.72 & 0.76 & 0.83 & 0.67 & 0.71 & 0.74 \\
 &  & MI & 0.29 & 0.23 & 0.95 & 0.63 & 0.78 & 0.34 & 0.56 \\
 &  & STTC & 0.27 & 0.17 & 0.97 & 0.61 & 0.76 & 0.27 & 0.51 \\
 &  & CD & 0.31 & 0.37 & 0.96 & 0.58 & 0.87 & 0.49 & 0.58 \\
 &  & HYP & 0.23 & 0.15 & 0.99 & 0.53 & 0.89 & 0.24 & 0.46 \\
\midrule
Fixed & PCA + LDA (grid 256--1024) & NORM & 0.39 & 0.72 & 0.74 & 0.88 & 0.57 & 0.71 & 0.73 \\
 &  & MI & 0.26 & 0.25 & 0.95 & 0.65 & 0.71 & 0.36 & 0.52 \\
 &  & STTC & 0.26 & 0.19 & 0.97 & 0.52 & 0.80 & 0.30 & 0.48 \\
 &  & CD & 0.23 & 0.34 & 0.98 & 0.63 & 0.84 & 0.48 & 0.59 \\
 &  & HYP & 0.13 & 0.29 & 0.98 & 0.59 & 0.90 & 0.41 & 0.50 \\
\midrule
Signal & 1D ResNet (raw signals) & NORM & 0.57 & 0.93 & 0.82 & 0.91 & 0.84 & 0.86 & 0.86 \\
 &  & MI & 0.35 & 0.67 & 0.95 & 0.76 & 0.91 & 0.74 & 0.76 \\
 &  & STTC & 0.35 & 0.76 & 0.93 & 0.85 & 0.88 & 0.76 & 0.77 \\
 &  & CD & 0.40 & 0.65 & 0.97 & 0.70 & 0.95 & 0.74 & 0.75 \\
 &  & HYP & 0.25 & 0.45 & 0.98 & 0.65 & 0.93 & 0.56 & 0.61 \\
\end{longtable}
\endgroup
{\footnotesize Thr.: mean validation-chosen threshold. Sens./Spec./F1 on the test fold at 0.5 and at the validation-chosen threshold.\par}
\medskip

\section{External per-label results}

\begingroup\scriptsize
\setlength{\tabcolsep}{3pt}
\begin{longtable}{lllrrrrrrrrr}
\caption{External AUROC per eligible superclass (mean over seeds). Left: Georgia, USA; right: China.}\label{tab:s_external_labels}\\
\toprule
Pages & Model & Runs & STTC & CD & HYP & Macro & MI & STTC & CD & HYP & Macro \\
\midrule
\endfirsthead
\toprule
Pages & Model & Runs & STTC & CD & HYP & Macro & MI & STTC & CD & HYP & Macro \\
\midrule
\endhead
\bottomrule
\endfoot
Adaptive & CNN-only & 5 & 0.678 & 0.693 & 0.585 & 0.652 & 0.686 & 0.750 & 0.727 & 0.704 & 0.717 \\
Adaptive & CNN-only, capacity-matched & 5 & 0.709 & 0.746 & 0.644 & 0.700 & 0.762 & 0.816 & 0.785 & 0.744 & 0.777 \\
Adaptive & CNN-LSTM & 5 & 0.723 & 0.750 & 0.697 & 0.723 & 0.789 & 0.825 & 0.764 & 0.819 & 0.799 \\
Adaptive & CNN + temporal conv. & 5 & 0.730 & 0.774 & 0.618 & 0.707 & 0.825 & 0.830 & 0.775 & 0.813 & 0.811 \\
Adaptive & CNN + small ViT & 5 & 0.727 & 0.764 & 0.699 & 0.730 & 0.832 & 0.825 & 0.771 & 0.833 & 0.815 \\
\midrule
Fixed & CNN-only & 5 & 0.696 & 0.712 & 0.713 & 0.707 & 0.744 & 0.775 & 0.736 & 0.817 & 0.768 \\
Fixed & CNN-only, capacity-matched$^\P$ & 5 & 0.711 & 0.751 & 0.739 & 0.733 & 0.808 & 0.813 & 0.769 & 0.847 & 0.809 \\
Fixed & CNN-LSTM & 5 & 0.731 & 0.760 & 0.724 & 0.738 & 0.815 & 0.826 & 0.773 & 0.854 & 0.817 \\
Fixed & CNN + temporal conv. & 5 & 0.732 & 0.768 & 0.713 & 0.738 & 0.841 & 0.834 & 0.776 & 0.846 & 0.824 \\
Fixed & CNN + small ViT & 5 & 0.728 & 0.767 & 0.717 & 0.737 & 0.841 & 0.827 & 0.761 & 0.846 & 0.819 \\
\midrule
Adaptive & PCA + LDA & 1 & 0.581 & 0.697 & 0.723 & 0.667 & 0.740 & 0.707 & 0.727 & 0.756 & 0.732 \\
Adaptive & PCA + linear SVM & 1 & 0.581 & 0.696 & 0.724 & 0.667 & 0.734 & 0.706 & 0.726 & 0.752 & 0.730 \\
\midrule
Fixed & PCA + LDA & 1 & 0.540 & 0.684 & 0.755 & 0.660 & 0.766 & 0.708 & 0.715 & 0.818 & 0.752 \\
\midrule
Signal & 1D ResNet (raw signals) & 5 & 0.747 & 0.853 & 0.761 & 0.787 & 0.875 & 0.827 & 0.849 & 0.871 & 0.855 \\
\end{longtable}
\endgroup
{\footnotesize Each run performed exactly one inference per cohort. NORM is not evaluable externally; MI is not evaluable in Georgia (7 positives). Signal seeds 45 and 46 were added under the 11 November 2026 amendment (Section~\ref{sec:s_amendments}) as a descriptive addition. $^\P$Additional analysis (amendment of 15 November 2026, Section~\ref{sec:s_review}): the capacity-matched head trained on fixed-gain pages, added after all other results were known. \par}
\medskip

\section{Gain policy: public resources and clipping at fixed gain}

\begingroup\scriptsize
\setlength{\tabcolsep}{3pt}
\begin{longtable}{p{3.6cm}p{3.6cm}p{1.6cm}p{3.6cm}p{2.2cm}}
\caption{How public ECG-image generators and datasets set the page gain (checked 11 November 2026).}\label{tab:s_survey}\\
\toprule
Resource & Gain when drawing & Amplitude-dependent? & Large deflections & Gain randomized? \\
\midrule
\endfirsthead
\toprule
Resource & Gain when drawing & Amplitude-dependent? & Large deflections & Gain randomized? \\
\midrule
\endhead
\bottomrule
\endfoot
ECG-Image-Kit \cite{shivashankara2024ecgimagekit} & fixed, 10 mm/mV; 1-mV pulse & no & no per-row clipping in code; traces may overlap adjacent rows & no (only pixel resolution) \\
PhysioNet Challenge 2024 images \cite{reyna2024challenge} & generated with ECG-Image-Kit & not stated & not stated & not stated \\
ECG-Image-Database \cite{reyna2024imagedatabase} & ECG-Image-Kit, printed; 0.5 mV per 5-mm square & not stated & not stated & not stated \\
Sangha et al.\ 2022 \cite{sangha2022images} & ecg-plot 0.2.8: fixed, $\approx$10.16 mm/mV at nominal size & no & no per-row clipping in code; traces may overlap adjacent rows & not stated \\
Khan et al.\ 2021 (real printouts) \cite{khan2021ecgimages} & not stated & not stated & not stated & not applicable \\
\end{longtable}
\endgroup
{\footnotesize Primary sources only (papers, code at a fixed commit, documentation); ``not stated'' means not found in those sources. ECG-Image-Kit: repository alphanumericslab/ecg-image-kit, commit 27b90f5, codes/ecg-image-generator/ecg\_plot.py (0.5 mV per 5-mm grid, page label 10 mm/mV). ecg-plot 0.2.8: the release current when Sangha et al.\ accessed it (repository dy1901/ecg\_plot, commit 3855e94). The project file docs/gain\_policy\_survey\_2026-11-11.md gives the URL and line or section for every entry.\par}
\medskip

\begingroup\footnotesize
\setlength{\tabcolsep}{3pt}
\begin{longtable}{rrrrrrr}
\caption{PTB-XL pages with at least one clipped sample if every page were drawn at one fixed gain (clipping at the edge of the 30-mm lead row).}\label{tab:s_clipping}\\
\toprule
Gain (mm/mV) & Pages & \% of 21{,}388 & Training & Validation & Test & HYP-positive \\
\midrule
\endfirsthead
\toprule
Gain (mm/mV) & Pages & \% of 21{,}388 & Training & Validation & Test & HYP-positive \\
\midrule
\endhead
\bottomrule
\endfoot
2.5 & 154 & 0.7 & 121 & 17 & 16 & 47 \\
5 & 2{,}541 & 11.9 & 2{,}031 & 251 & 259 & 1{,}033 \\
10 & 16{,}494 & 77.1 & 13{,}168 & 1{,}657 & 1{,}669 & 2{,}516 \\
\end{longtable}
\endgroup
{\footnotesize Definition used in this table: a page is counted when at least one sample of at least one lead falls outside that lead's 30-mm row ($\pm$15 mm around the row centre); the fixed-gain renderer clips such samples to the row edge. Numerator: pages with $\geq$1 clipped sample; denominator: all pages of the column (21{,}388 in total; 17{,}084 training, 2{,}146 validation, 2{,}158 test; 2{,}649 HYP-positive). Stage: page rendering, before any model sees the page; no page was excluded. The 2.5 mm/mV counts reproduce the clipping recorded in the fixed-gain dataset for every ECG. The 121 training pages of the first row are a subset of the 274 training pages that exceed the narrower $\pm$12-mm band used to choose the gain; Supplementary Table~\ref{tab:s_clipping_defs} gives both counts side by side, for PTB-XL and for both external cohorts. Computed by src/count\_gain\_clipping.py with the unchanged renderer.\par}
\medskip

\section{Sensitivity analyses of 11 November 2026: a pretrained ResNet-18 and a longer training budget}

\begingroup\scriptsize
\setlength{\tabcolsep}{3pt}
\begin{longtable}{lllrrrrrrrr}
\caption{ImageNet-pretrained ResNet-18 (sensitivity analysis, amendment of 11 November 2026): every run.}\label{tab:s_resnet_runs}\\
\toprule
Cohort & Pages & Seed & Macro & NORM & MI & STTC & CD & HYP & $\Delta$ & Epoch \\
\midrule
\endfirsthead
\toprule
Cohort & Pages & Seed & Macro & NORM & MI & STTC & CD & HYP & $\Delta$ & Epoch \\
\midrule
\endhead
\bottomrule
\endfoot
PTB-XL test & Adaptive & 42 & 0.920 & 0.941 & 0.930 & 0.929 & 0.917 & 0.884 & -- & 8 \\
PTB-XL test & Adaptive & 43 & 0.920 & 0.940 & 0.927 & 0.932 & 0.919 & 0.883 & -- & 7 \\
PTB-XL test & Adaptive & 44 & 0.922 & 0.943 & 0.924 & 0.932 & 0.921 & 0.889 & -- & 8 \\
PTB-XL test & Fixed & 42 & 0.929 & 0.950 & 0.931 & 0.935 & 0.922 & 0.908 & -- & 7 \\
PTB-XL test & Fixed & 43 & 0.931 & 0.949 & 0.934 & 0.936 & 0.925 & 0.911 & -- & 10 \\
PTB-XL test & Fixed & 44 & 0.927 & 0.946 & 0.925 & 0.935 & 0.924 & 0.904 & -- & 9 \\
\midrule
Georgia, USA & Adaptive & 42 & 0.772 & -- & 0.775 & 0.749 & 0.850 & 0.718 & $-$0.138 & -- \\
Georgia, USA & Adaptive & 43 & 0.774 & -- & 0.725 & 0.752 & 0.839 & 0.731 & $-$0.137 & -- \\
Georgia, USA & Adaptive & 44 & 0.774 & -- & 0.738 & 0.755 & 0.850 & 0.718 & $-$0.140 & -- \\
Georgia, USA & Fixed & 42 & 0.780 & -- & 0.759 & 0.752 & 0.858 & 0.729 & $-$0.142 & -- \\
Georgia, USA & Fixed & 43 & 0.782 & -- & 0.784 & 0.754 & 0.845 & 0.746 & $-$0.142 & -- \\
Georgia, USA & Fixed & 44 & 0.774 & -- & 0.754 & 0.745 & 0.855 & 0.724 & $-$0.147 & -- \\
\midrule
China & Adaptive & 42 & 0.850 & -- & 0.898 & 0.812 & 0.844 & 0.847 & $-$0.065 & -- \\
China & Adaptive & 43 & 0.845 & -- & 0.875 & 0.835 & 0.853 & 0.816 & $-$0.070 & -- \\
China & Adaptive & 44 & 0.849 & -- & 0.887 & 0.834 & 0.842 & 0.835 & $-$0.067 & -- \\
China & Fixed & 42 & 0.851 & -- & 0.898 & 0.825 & 0.837 & 0.842 & $-$0.073 & -- \\
China & Fixed & 43 & 0.855 & -- & 0.902 & 0.829 & 0.843 & 0.845 & $-$0.072 & -- \\
China & Fixed & 44 & 0.853 & -- & 0.913 & 0.815 & 0.837 & 0.847 & $-$0.069 & -- \\
\end{longtable}
\endgroup
{\footnotesize Macro AUROC over all five labels on PTB-XL and over the eligible labels externally (Georgia: STTC, CD, HYP; China: MI, STTC, CD, HYP). Per-label columns are AUROC (labels not eligible externally are shown for completeness). $\Delta$: external minus PTB-XL test macro AUROC on the same labels. Epoch: selected epoch (of at most 20). Every reloaded model reproduced its saved PTB-XL validation probabilities exactly before its single external inference.\par}
\medskip

\begingroup\footnotesize
\setlength{\tabcolsep}{3pt}
\begin{longtable}{llcr}
\caption{ImageNet-pretrained ResNet-18: fixed-gain minus adaptive-gain AUROC [95\% CI] and Holm-adjusted $p$ (seeds 42--44).}\label{tab:s_resnet_fixed}\\
\toprule
Cohort & Metric & Difference & $p$ \\
\midrule
\endfirsthead
\toprule
Cohort & Metric & Difference & $p$ \\
\midrule
\endhead
\bottomrule
\endfoot
PTB-XL test & macro AUROC & +0.008 [+0.003, +0.013] & 0.001 \\
PTB-XL test & HYP AUROC & +0.022 [+0.009, +0.035] & 0.001 \\
Georgia, USA & macro AUROC & +0.005 [$-$0.000, +0.010] & 0.175 \\
Georgia, USA & HYP AUROC & +0.011 [+0.002, +0.019] & 0.076 \\
China & macro AUROC & +0.004 [$-$0.002, +0.011] & 0.302 \\
China & HYP AUROC & +0.012 [$-$0.007, +0.030] & 0.302 \\
\end{longtable}
\endgroup
{\footnotesize PTB-XL: hierarchical seed $\times$ patient bootstrap, Holm over 2 comparisons. External cohorts: seed $\times$ record bootstrap, Holm over the 4 cohort $\times$ metric comparisons. 10{,}000 resamples, seed 2024. $p$ printed as ``$\le$'' is at the resolution limit of the bootstrap: no replicate crossed zero, and no smaller value is claimed (Section~\ref{sec:s_review}). Prespecified in the amendment of 11 November 2026, after the frozen external plan; not part of the confirmatory hypotheses H1 and H2.\par}
\medskip

\begingroup\scriptsize
\setlength{\tabcolsep}{3pt}
\begin{longtable}{lrcrll}
\caption{Training-budget sensitivity (amendment of 11 November 2026): 20 versus at most 40 epochs, adaptive-gain pages, seeds 42--44.}\label{tab:s_long40}\\
\toprule
Model & Max. epochs & Test macro AUROC & Validation & Selected epochs & Epochs run \\
\midrule
\endfirsthead
\toprule
Model & Max. epochs & Test macro AUROC & Validation & Selected epochs & Epochs run \\
\midrule
\endhead
\bottomrule
\endfoot
CNN-only & 20 & 0.829 $\pm$ 0.005 & 0.841 & 19 20 19 & 20 20 20 \\
CNN-only & 40 & 0.845 $\pm$ 0.003 & 0.858 & 33 37 36 & 37 40 40 \\
CNN-only, capacity-m. & 20 & 0.865 $\pm$ 0.003 & 0.876 & 20 20 18 & 20 20 20 \\
CNN-only, capacity-m. & 40 & 0.884 $\pm$ 0.001 & 0.892 & 37 31 35 & 40 35 39 \\
CNN-LSTM & 20 & 0.879 $\pm$ 0.008 & 0.893 & 15 16 18 & 19 20 20 \\
CNN-LSTM & 40 & 0.891 $\pm$ 0.003 & 0.903 & 27 26 28 & 31 30 32 \\
CNN + temporal conv. & 20 & 0.879 $\pm$ 0.001 & 0.887 & 16 19 20 & 20 20 20 \\
CNN + temporal conv. & 40 & 0.891 $\pm$ 0.002 & 0.899 & 29 35 30 & 33 39 34 \\
CNN + small ViT & 20 & 0.888 $\pm$ 0.002 & 0.898 & 19 18 13 & 20 20 17 \\
CNN + small ViT & 40 & 0.891 $\pm$ 0.006 & 0.900 & 15 30 21 & 19 34 25 \\
\end{longtable}
\endgroup
{\footnotesize Identical settings except the epoch budget (the cosine schedule spans the budget; early-stopping patience 4 in both). Seed 42 of the 20-epoch CNN-only, CNN-LSTM, temporal-convolution and ViT runs was trained on a CPU; all 40-epoch runs on the GPU.\par}
\medskip

\begingroup\scriptsize
\setlength{\tabcolsep}{3pt}
\begin{longtable}{llcr}
\caption{Training-budget sensitivity: 40 minus 20 epochs per model, and the ten head pairs with each budget (PTB-XL test, seeds 42--44).}\label{tab:s_long40_pairs}\\
\toprule
Family & Comparison & Difference [95\% CI] & $p$ \\
\midrule
\endfirsthead
\toprule
Family & Comparison & Difference [95\% CI] & $p$ \\
\midrule
\endhead
\bottomrule
\endfoot
40 - 20 epochs & CNN-only & +0.016 [+0.011, +0.020] & $\le$0.001 \\
40 - 20 epochs & CNN-only, capacity-m. & +0.019 [+0.015, +0.023] & $\le$0.001 \\
40 - 20 epochs & CNN-LSTM & +0.012 [+0.005, +0.021] & $\le$0.001 \\
40 - 20 epochs & CNN + temporal conv. & +0.012 [+0.008, +0.016] & $\le$0.001 \\
40 - 20 epochs & CNN + small ViT & +0.003 [$-$0.003, +0.010] & 0.312 \\
\midrule
head pairs, 40 epochs & CNN-only, capacity-m. $-$ CNN-only & +0.039 [+0.032, +0.045] & $\le$0.002 \\
head pairs, 40 epochs & CNN-LSTM $-$ CNN-only & +0.046 [+0.038, +0.055] & $\le$0.002 \\
head pairs, 40 epochs & CNN + temporal conv. $-$ CNN-only & +0.047 [+0.039, +0.053] & $\le$0.002 \\
head pairs, 40 epochs & CNN + small ViT $-$ CNN-only & +0.046 [+0.038, +0.055] & $\le$0.002 \\
head pairs, 40 epochs & CNN-LSTM $-$ CNN-only, capacity-m. & +0.007 [+0.001, +0.013] & 0.110 \\
head pairs, 40 epochs & CNN + temporal conv. $-$ CNN-only, capacity-m. & +0.008 [+0.003, +0.013] & 0.026 \\
head pairs, 40 epochs & CNN + small ViT $-$ CNN-only, capacity-m. & +0.008 [+0.001, +0.015] & 0.110 \\
head pairs, 40 epochs & CNN + temporal conv. $-$ CNN-LSTM & +0.000 [$-$0.007, +0.008] & 1.000 \\
head pairs, 40 epochs & CNN + small ViT $-$ CNN-LSTM & +0.000 [$-$0.008, +0.010] & 1.000 \\
head pairs, 40 epochs & CNN + small ViT $-$ CNN + temporal conv. & $-$0.000 [$-$0.006, +0.006] & 1.000 \\
\midrule
head pairs, 20 epochs (same seeds) & CNN-only, capacity-m. $-$ CNN-only & +0.036 [+0.028, +0.043] & $\le$0.002 \\
head pairs, 20 epochs (same seeds) & CNN-LSTM $-$ CNN-only & +0.050 [+0.036, +0.063] & $\le$0.002 \\
head pairs, 20 epochs (same seeds) & CNN + temporal conv. $-$ CNN-only & +0.050 [+0.043, +0.058] & $\le$0.002 \\
head pairs, 20 epochs (same seeds) & CNN + small ViT $-$ CNN-only & +0.059 [+0.051, +0.067] & $\le$0.002 \\
head pairs, 20 epochs (same seeds) & CNN-LSTM $-$ CNN-only, capacity-m. & +0.014 [+0.004, +0.024] & 0.024 \\
head pairs, 20 epochs (same seeds) & CNN + temporal conv. $-$ CNN-only, capacity-m. & +0.014 [+0.009, +0.020] & $\le$0.002 \\
head pairs, 20 epochs (same seeds) & CNN + small ViT $-$ CNN-only, capacity-m. & +0.023 [+0.017, +0.029] & $\le$0.002 \\
head pairs, 20 epochs (same seeds) & CNN + temporal conv. $-$ CNN-LSTM & +0.000 [$-$0.009, +0.010] & 0.943 \\
head pairs, 20 epochs (same seeds) & CNN + small ViT $-$ CNN-LSTM & +0.009 [$-$0.001, +0.021] & 0.159 \\
head pairs, 20 epochs (same seeds) & CNN + small ViT $-$ CNN + temporal conv. & +0.009 [+0.004, +0.014] & 0.006 \\
\end{longtable}
\endgroup
{\footnotesize Difference in mean-over-seeds test macro AUROC; hierarchical seed $\times$ patient bootstrap (10{,}000 resamples), Holm within each family (5 or 10 comparisons). $p$ printed as ``$\le$'' is at the resolution limit of the bootstrap: no replicate crossed zero, and no smaller value is claimed (Section~\ref{sec:s_review}). \par}
\medskip

\begingroup\scriptsize
\setlength{\tabcolsep}{3pt}
\begin{longtable}{lllcr}
\caption{Fixed-gain minus adaptive-gain training with a 40-epoch budget, fixed-gain 40 minus 20 epochs, and the 20-epoch fixed-minus-adaptive differences for the same seeds (PTB-XL test, seeds 42--44).}\label{tab:s_fg_long40}\\
\toprule
Family & Model & AUROC & Difference [95\% CI] & $p$ \\
\midrule
\endfirsthead
\toprule
Family & Model & AUROC & Difference [95\% CI] & $p$ \\
\midrule
\endhead
\bottomrule
\endfoot
fixed - adaptive, 40 epochs & CNN-only & macro & +0.032 [+0.023, +0.041] & $\le$0.002 \\
fixed - adaptive, 40 epochs & CNN-only & HYP & +0.060 [+0.036, +0.087] & $\le$0.002 \\
fixed - adaptive, 40 epochs & CNN-LSTM & macro & +0.002 [$-$0.005, +0.011] & 1.000 \\
fixed - adaptive, 40 epochs & CNN-LSTM & HYP & +0.000 [$-$0.013, +0.015] & 1.000 \\
fixed - adaptive, 40 epochs & CNN + temporal conv. & macro & +0.005 [$-$0.001, +0.012] & 0.527 \\
fixed - adaptive, 40 epochs & CNN + temporal conv. & HYP & +0.023 [+0.010, +0.036] & 0.004 \\
fixed - adaptive, 40 epochs & CNN + small ViT & macro & +0.004 [$-$0.007, +0.013] & 1.000 \\
fixed - adaptive, 40 epochs & CNN + small ViT & HYP & +0.005 [$-$0.013, +0.024] & 1.000 \\
\midrule
fixed gain, 40 - 20 epochs & CNN-only & macro & +0.016 [+0.013, +0.020] & $\le$0.001 \\
fixed gain, 40 - 20 epochs & CNN-LSTM & macro & +0.008 [+0.003, +0.013] & 0.002 \\
fixed gain, 40 - 20 epochs & CNN + temporal conv. & macro & +0.002 [$-$0.003, +0.007] & 0.456 \\
fixed gain, 40 - 20 epochs & CNN + small ViT & macro & +0.003 [$-$0.001, +0.009] & 0.346 \\
\midrule
fixed - adaptive, 20 epochs (same seeds) & CNN-only & macro & +0.032 [+0.019, +0.043] & $\le$0.002 \\
fixed - adaptive, 20 epochs (same seeds) & CNN-only & HYP & +0.073 [+0.045, +0.104] & $\le$0.002 \\
fixed - adaptive, 20 epochs (same seeds) & CNN-LSTM & macro & +0.007 [$-$0.001, +0.015] & 0.361 \\
fixed - adaptive, 20 epochs (same seeds) & CNN-LSTM & HYP & +0.011 [$-$0.006, +0.029] & 0.649 \\
fixed - adaptive, 20 epochs (same seeds) & CNN + temporal conv. & macro & +0.015 [+0.010, +0.020] & $\le$0.002 \\
fixed - adaptive, 20 epochs (same seeds) & CNN + temporal conv. & HYP & +0.040 [+0.026, +0.056] & $\le$0.002 \\
fixed - adaptive, 20 epochs (same seeds) & CNN + small ViT & macro & +0.003 [$-$0.006, +0.012] & 0.649 \\
fixed - adaptive, 20 epochs (same seeds) & CNN + small ViT & HYP & +0.008 [$-$0.007, +0.023] & 0.649 \\
\end{longtable}
\endgroup
{\footnotesize Fixed-gain amendment of 11 November 2026 (evening): the four fixed-gain networks retrained with at most 40 epochs, otherwise identical, seeds 42--44, all on the GPU. Selected epochs (epochs run): CNN-only 40 39 40 (40 40 40); CNN-LSTM 24 27 26 (28 31 30); CNN + temporal conv. 11 27 20 (15 31 24); CNN + small ViT 19 27 14 (23 31 18). Hierarchical seed $\times$ patient bootstrap on the PTB-XL test fold (10{,}000 resamples), Holm within each family (8, 4 and 8 comparisons). $p$ printed as ``$\le$'' is at the resolution limit of the bootstrap: no replicate crossed zero, and no smaller value is claimed (Section~\ref{sec:s_review}). No external inference was run for these models. The statement in the amendment text below that the capacity-matched head had no fixed-gain runs was true on 11 November; those runs were added on 15 November (Section~\ref{sec:s_review}).\par}
\medskip

\section{Additional analyses of 15 November 2026 (review by the senior author)}

\label{sec:s_review}Everything in this section was added on 15 November 2026 in response to a review of the manuscript by the senior author, after all earlier results were known. Each part follows a dated protocol that was written and committed before the corresponding numbers existed (reproduced in Section~\ref{sec:s_amendments}). The analyses are labelled \textbf{Additional analysis} here and in the main text and are kept apart from the prospectively specified analyses; none of them changes a completed run, a saved prediction or the frozen external plan. Source: results/review\_20261115/, generated by src/review\_20261115.py.

\subsection{Capacity-matched head on fixed-gain pages}

\begingroup\scriptsize
\setlength{\tabcolsep}{3pt}
\begin{longtable}{lllrrrrrrrrr}
\caption{Additional analysis: capacity-matched CNN-only head trained on fixed-gain (2.5 mm/mV) pages, every run (PTB-XL; test fold exploratory).}\label{tab:s_wide_runs}\\
\toprule
Seed & Device & Selected epoch & Val. & Test & AP & F1 & NORM & MI & STTC & CD & HYP \\
\midrule
\endfirsthead
\toprule
Seed & Device & Selected epoch & Val. & Test & AP & F1 & NORM & MI & STTC & CD & HYP \\
\midrule
\endhead
\bottomrule
\endfoot
42 & GPU & 18 of 20 & 0.892 & 0.882 & 0.741 & 0.643 & 0.916 & 0.849 & 0.908 & 0.853 & 0.886 \\
43 & GPU & 18 of 20 & 0.888 & 0.875 & 0.725 & 0.628 & 0.911 & 0.838 & 0.905 & 0.847 & 0.872 \\
44 & GPU & 17 of 20 & 0.889 & 0.880 & 0.733 & 0.640 & 0.916 & 0.843 & 0.911 & 0.856 & 0.874 \\
45 & GPU & 20 of 20 & 0.897 & 0.886 & 0.744 & 0.645 & 0.917 & 0.872 & 0.909 & 0.861 & 0.872 \\
46 & GPU & 20 of 20 & 0.899 & 0.890 & 0.754 & 0.653 & 0.925 & 0.863 & 0.919 & 0.863 & 0.878 \\
\midrule
Mean &  &  & 0.893 & 0.883 & 0.739 & 0.642 & 0.917 & 0.853 & 0.910 & 0.856 & 0.876 \\
SD &  &  & 0.005 & 0.006 & 0.011 & 0.009 & 0.005 & 0.014 & 0.005 & 0.006 & 0.006 \\
\end{longtable}
\endgroup
{\footnotesize Protocol: amendment of 15 November 2026 (capacity-matched head on fixed-gain pages), frozen before the first run. Config experiments/configs/ptbxl\_v2\_fg\_cnn\_pool\_wide.json: the head's adaptive-gain config with only the dataset fields changed (dataset ptbxl\_images\_fixedgain2p5\_v1); identical training block, the same five seeds, the same validation-based selection, threshold 0.5 and one test inference per run. Val./Test: macro AUROC; AP: macro average precision; F1: macro F1 at 0.5; label columns: test AUROC. Run identifiers: ptbxl\_v2\_fg\_cnn\_pool\_wide\_seed42 to seed46 (results/ptbxl\_v2\_fixedgain/), run on 2026-11-15 (UTC).\par}
\medskip

\begingroup\footnotesize
\setlength{\tabcolsep}{3pt}
\begin{longtable}{llccc}
\caption{Additional analysis: capacity-matched head, fixed-gain minus adaptive-gain training: difference in mean-over-seeds AUROC [95\% CI] and Holm-adjusted $p$ (seeds 42--46).}\label{tab:s_wide_fixed}\\
\toprule
Cohort & AUROC & Difference & $p$ (family of 2) & $p$ (family of 10) \\
\midrule
\endfirsthead
\toprule
Cohort & AUROC & Difference & $p$ (family of 2) & $p$ (family of 10) \\
\midrule
\endhead
\bottomrule
\endfoot
PTB-XL test & macro & +0.017 [+0.010, +0.024] & $\le$0.001 & $\le$0.002 \\
PTB-XL test & HYP & +0.043 [+0.019, +0.069] & $\le$0.001 & $\le$0.002 \\
Georgia, USA & macro & +0.034 [+0.020, +0.048] & $\le$0.001 & $\le$0.002 \\
Georgia, USA & HYP & +0.095 [+0.047, +0.141] & $\le$0.001 & $\le$0.002 \\
China & macro & +0.033 [+0.021, +0.043] & $\le$0.001 & $\le$0.002 \\
China & HYP & +0.103 [+0.075, +0.131] & $\le$0.001 & $\le$0.002 \\
\end{longtable}
\endgroup
{\footnotesize Macro AUROC over the eligible labels of each cohort; HYP: hypertrophy AUROC. Hierarchical bootstrap (seeds, then patients on PTB-XL or records externally; 10{,}000 resamples). Family of 2: the two metrics of this head within the cohort (prespecified primary). Family of 10: this head together with the four heads of the main fixed-gain analysis (five heads $\times$ two metrics), as a sensitivity analysis. All ten runs of this head were GPU-trained. External cohorts: one inference per run, outside the confirmatory hypotheses H1 and H2. $p$ printed as ``$\le$'' is at the resolution limit of the bootstrap: no replicate crossed zero, and no smaller value is claimed (Section~\ref{sec:s_review}). \par}
\medskip

\begingroup\scriptsize
\setlength{\tabcolsep}{3pt}
\begin{longtable}{lcccccc}
\caption{Additional analysis: head pairs on fixed-gain pages: difference in mean-over-seeds macro AUROC [95\% CI] and Holm-adjusted $p$ (seeds 42--46).}\label{tab:s_fixed_pairs}\\
\toprule
Comparison & \multicolumn{2}{c}{PTB-XL test} & \multicolumn{2}{c}{Georgia, USA} & \multicolumn{2}{c}{China} \\
\midrule
\endfirsthead
\toprule
Comparison & \multicolumn{2}{c}{PTB-XL test} & \multicolumn{2}{c}{Georgia, USA} & \multicolumn{2}{c}{China} \\
\midrule
\endhead
\bottomrule
\endfoot
CNN-only, capacity-m. $-$ CNN-only & +0.022 [+0.014, +0.030] & $\le$0.002 & +0.026 [+0.020, +0.033] & $\le$0.002 & +0.042 [+0.027, +0.054] & $\le$0.002 \\
CNN-LSTM $-$ CNN-only & +0.029 [+0.021, +0.036] & $\le$0.002 & +0.031 [+0.023, +0.039] & $\le$0.002 & +0.049 [+0.038, +0.060] & $\le$0.002 \\
CNN + temporal conv. $-$ CNN-only & +0.035 [+0.028, +0.042] & $\le$0.002 & +0.031 [+0.022, +0.040] & $\le$0.002 & +0.057 [+0.040, +0.072] & $\le$0.002 \\
CNN + small ViT $-$ CNN-only & +0.031 [+0.020, +0.042] & $\le$0.002 & +0.030 [+0.021, +0.040] & $\le$0.002 & +0.051 [+0.039, +0.064] & $\le$0.002 \\
CNN-LSTM $-$ CNN-only, capacity-m. & +0.006 [+0.001, +0.011] & 0.081 & +0.005 [$-$0.002, +0.010] & 0.632 & +0.008 [+0.001, +0.014] & 0.110 \\
CNN + temporal conv. $-$ CNN-only, capacity-m. & +0.012 [+0.007, +0.018] & $\le$0.002 & +0.004 [$-$0.000, +0.009] & 0.334 & +0.015 [+0.007, +0.022] & 0.004 \\
CNN + small ViT $-$ CNN-only, capacity-m. & +0.008 [+0.001, +0.015] & 0.081 & +0.004 [$-$0.002, +0.010] & 0.737 & +0.009 [+0.003, +0.016] & 0.029 \\
CNN + temporal conv. $-$ CNN-LSTM & +0.006 [+0.001, +0.011] & 0.081 & $-$0.000 [$-$0.007, +0.005] & 1.000 & +0.007 [$-$0.004, +0.018] & 0.595 \\
CNN + small ViT $-$ CNN-LSTM & +0.002 [$-$0.005, +0.010] & 0.571 & $-$0.001 [$-$0.006, +0.004] & 1.000 & +0.002 [$-$0.005, +0.009] & 0.626 \\
CNN + small ViT $-$ CNN + temporal conv. & $-$0.004 [$-$0.010, +0.002] & 0.439 & $-$0.001 [$-$0.005, +0.003] & 1.000 & $-$0.006 [$-$0.014, +0.003] & 0.595 \\
\end{longtable}
\endgroup
{\footnotesize The same ten pairs as Table~2 of the main text, with every model trained and tested on fixed-gain pages. Holm over the 10 pairs within each cohort; 10{,}000 resamples. $p$ printed as ``$\le$'' is at the resolution limit of the bootstrap: no replicate crossed zero, and no smaller value is claimed (Section~\ref{sec:s_review}). \par}
\medskip

\subsection{PCA + LDA under fixed gain, and the gap to the signal reference}

\begingroup\footnotesize
\setlength{\tabcolsep}{3pt}
\begin{longtable}{llcc}
\caption{Additional analysis: PCA + LDA, fixed-gain minus adaptive-gain pages: difference in AUROC [95\% CI] and Holm-adjusted $p$.}\label{tab:s_lda_fixed}\\
\toprule
Cohort & AUROC & Difference & $p$ \\
\midrule
\endfirsthead
\toprule
Cohort & AUROC & Difference & $p$ \\
\midrule
\endhead
\bottomrule
\endfoot
PTB-XL test & macro & $-$0.013 [$-$0.023, $-$0.003] & 0.017 \\
PTB-XL test & HYP & +0.017 [$-$0.005, +0.039] & 0.119 \\
Georgia, USA & macro & $-$0.007 [$-$0.013, $-$0.002] & 0.016 \\
Georgia, USA & HYP & +0.032 [+0.022, +0.042] & $\le$0.001 \\
China & macro & +0.019 [+0.012, +0.027] & $\le$0.001 \\
China & HYP & +0.062 [+0.049, +0.075] & $\le$0.001 \\
\end{longtable}
\endgroup
{\footnotesize The two single fits reported in the main text (no training seed, so the interval reflects test sampling only: patients resampled on PTB-XL, records externally; 10{,}000 resamples). Holm over the 2 metrics within each cohort. Post hoc. $p$ printed as ``$\le$'' is at the resolution limit of the bootstrap: no replicate crossed zero, and no smaller value is claimed (Section~\ref{sec:s_review}). \par}
\medskip

\begingroup\scriptsize
\setlength{\tabcolsep}{3pt}
\begin{longtable}{lllcccccc}
\caption{Additional analysis: image model minus signal reference: difference in mean-over-seeds macro AUROC [95\% CI] and Holm-adjusted $p$.}\label{tab:s_signal_gap}\\
\toprule
Model & Pages & Seeds & \multicolumn{2}{c}{PTB-XL test} & \multicolumn{2}{c}{Georgia, USA} & \multicolumn{2}{c}{China} \\
\midrule
\endfirsthead
\toprule
Model & Pages & Seeds & \multicolumn{2}{c}{PTB-XL test} & \multicolumn{2}{c}{Georgia, USA} & \multicolumn{2}{c}{China} \\
\midrule
\endhead
\bottomrule
\endfoot
CNN-only & Adaptive & 5 & $-$0.097 [$-$0.107, $-$0.088] & $\le$0.003 & $-$0.135 [$-$0.143, $-$0.127] & $\le$0.003 & $-$0.139 [$-$0.153, $-$0.126] & $\le$0.003 \\
CNN-only & Fixed & 5 & $-$0.067 [$-$0.076, $-$0.059] & $\le$0.003 & $-$0.080 [$-$0.087, $-$0.073] & $\le$0.003 & $-$0.088 [$-$0.100, $-$0.075] & $\le$0.003 \\
CNN-only, capacity-matched & Adaptive & 5 & $-$0.062 [$-$0.070, $-$0.055] & $\le$0.003 & $-$0.087 [$-$0.102, $-$0.071] & $\le$0.003 & $-$0.079 [$-$0.092, $-$0.067] & $\le$0.003 \\
CNN-only, capacity-matched & Fixed & 5 & $-$0.045 [$-$0.053, $-$0.038] & $\le$0.003 & $-$0.053 [$-$0.058, $-$0.048] & $\le$0.003 & $-$0.046 [$-$0.055, $-$0.036] & $\le$0.003 \\
CNN-LSTM & Adaptive & 5 & $-$0.051 [$-$0.059, $-$0.043] & $\le$0.003 & $-$0.064 [$-$0.072, $-$0.056] & $\le$0.003 & $-$0.056 [$-$0.064, $-$0.049] & $\le$0.003 \\
CNN-LSTM & Fixed & 5 & $-$0.039 [$-$0.046, $-$0.032] & $\le$0.003 & $-$0.048 [$-$0.055, $-$0.043] & $\le$0.003 & $-$0.038 [$-$0.048, $-$0.028] & $\le$0.003 \\
CNN + temporal conv. & Adaptive & 5 & $-$0.048 [$-$0.055, $-$0.042] & $\le$0.003 & $-$0.080 [$-$0.086, $-$0.073] & $\le$0.003 & $-$0.045 [$-$0.054, $-$0.035] & $\le$0.003 \\
CNN + temporal conv. & Fixed & 5 & $-$0.033 [$-$0.038, $-$0.028] & $\le$0.003 & $-$0.049 [$-$0.054, $-$0.043] & $\le$0.003 & $-$0.031 [$-$0.039, $-$0.022] & $\le$0.003 \\
CNN + small ViT & Adaptive & 5 & $-$0.042 [$-$0.049, $-$0.036] & $\le$0.003 & $-$0.057 [$-$0.064, $-$0.049] & $\le$0.003 & $-$0.040 [$-$0.047, $-$0.033] & $\le$0.003 \\
CNN + small ViT & Fixed & 5 & $-$0.036 [$-$0.044, $-$0.029] & $\le$0.003 & $-$0.049 [$-$0.056, $-$0.043] & $\le$0.003 & $-$0.037 [$-$0.044, $-$0.029] & $\le$0.003 \\
\midrule
ResNet-18, ImageNet & Adaptive & 3 & $-$0.008 [$-$0.012, $-$0.003] & 0.002 & $-$0.015 [$-$0.019, $-$0.010] & $\le$0.003 & $-$0.007 [$-$0.015, +0.000] & 0.104 \\
ResNet-18, ImageNet & Fixed & 3 & +0.001 [$-$0.004, +0.005] & 0.801 & $-$0.010 [$-$0.016, $-$0.004] & 0.002 & $-$0.003 [$-$0.010, +0.004] & 0.376 \\
\end{longtable}
\endgroup
{\footnotesize Negative values: the image model is below the 1D ResNet on raw signals. Compact heads: seeds 42--46; ResNet-18: seeds 42--44 (the signal reference restricted to the same seeds). Hierarchical bootstrap, 10{,}000 resamples, Holm over the 12 comparisons within each cohort. Post hoc; seeds of different architectures are paired by number only. An interval that includes zero means the two models were not distinguishable in this sample; it does not show that they are equivalent. $p$ printed as ``$\le$'' is at the resolution limit of the bootstrap: no replicate crossed zero, and no smaller value is claimed (Section~\ref{sec:s_review}). \par}
\medskip

\subsection{Clipping: two boundaries and the external cohorts}

\begingroup\scriptsize
\setlength{\tabcolsep}{3pt}
\begin{longtable}{llrrrrrrrr}
\caption{Additional analysis: the two amplitude boundaries at 2.5 mm/mV, for PTB-XL and both external cohorts.}\label{tab:s_clipping_defs}\\
\toprule
Cohort & Records & Pages & Outside band & \% & Clipped & \% & HYP+ & HYP+ outside band & HYP+ clipped \\
\midrule
\endfirsthead
\toprule
Cohort & Records & Pages & Outside band & \% & Clipped & \% & HYP+ & HYP+ outside band & HYP+ clipped \\
\midrule
\endhead
\bottomrule
\endfoot
PTB-XL & training (folds 1--8) & 17{,}084 & 274 & 1.60 & 121 & 0.71 & 2{,}119 & 104 & 37 \\
PTB-XL & validation (fold 9) & 2{,}146 & 40 & 1.86 & 17 & 0.79 & 268 & 13 & 2 \\
PTB-XL & test (fold 10) & 2{,}158 & 36 & 1.67 & 16 & 0.74 & 262 & 16 & 8 \\
PTB-XL & all & 21{,}388 & 350 & 1.64 & 154 & 0.72 & 2{,}649 & 133 & 47 \\
\midrule
Georgia, USA & included (confirmatory cohort) & 6{,}047 & 78 & 1.29 & 28 & 0.46 & 2{,}088 & 47 & 14 \\
Georgia, USA & no mapped superclass (sensitivity only) & 4{,}245 & 21 & 0.49 & 14 & 0.33 & 0 & 0 & 0 \\
Georgia, USA & all rendered records & 10{,}292 & 99 & 0.96 & 42 & 0.41 & 2{,}088 & 47 & 14 \\
\midrule
China & included (confirmatory cohort) & 9{,}004 & 795 & 8.83 & 383 & 4.25 & 770 & 225 & 75 \\
\end{longtable}
\endgroup
{\footnotesize \emph{Outside band}: at 2.5 mm/mV at least one sample lies beyond the $\pm$12-mm band that the original layout allows around a row centre. This is the criterion by which 2.5 mm/mV was chosen (at least 98\% of training pages must fit: 16{,}810 of 17{,}084 fit, 274 do not) and it is identical to the set of ECGs that the adaptive renderer drew below 2.5 mm/mV. Nothing is clipped at this boundary. \emph{Clipped}: at least one sample lies beyond the edge of the 30-mm row ($\pm$15 mm) and was clipped to the row edge by the fixed-gain renderer. Every clipped page is also outside the band (checked), so the 121 clipped training pages are a subset of the 274. Numerators: pages meeting the definition; denominators: the pages of the row (column Pages). Stage: page rendering; no page was excluded for clipping in any cohort. HYP+: pages with the hypertrophy label (external cohorts: the mapped label). ``No mapped superclass'' rows are the external records used only in the label-mapping sensitivity analysis. Source: the fixed-gain cache manifests.\par}
\medskip

\subsection{Label mapping: clinical names, excluded records and sensitivity}

\paragraph{Validation of the label mapping without an expert.} No cardiologist reviewed the mapping; this remains a limitation. Three checks were possible. (i) \emph{Names.} Every code has a clinical name in the Challenge organizers' code tables (Supplementary Table~\ref{tab:s_mapping}), and the name agrees in meaning with the SNOMED CT concept for all assigned codes except anterior myocardial infarction (54329005), whose SNOMED CT concept is \emph{acute} myocardial infarction of the anterior wall. (ii) \emph{Clinical meaning against the PTB-XL superclass definitions.} The assignments of the MI, STTC and CD codes agree with PTB-XL's own statement table: ischemic and nonspecific ST-T changes and a prolonged QT interval belong to STTC, and bundle-branch, fascicular and atrioventricular blocks and the Wolff-Parkinson-White pattern to CD. For HYP, left atrial enlargement and right atrial hypertrophy are assigned because PTB-XL places atrial overload/enlargement in its hypertrophy superclass. One assignment has no clinical basis: right ventricular hypertrophy (89792004) is assigned to CD as well as HYP only because 94\% of the 126 PTB-XL records with it also have a conduction disturbance (co-occurrence); PTB-XL itself classifies right ventricular hypertrophy as HYP only. Variant V1 removes this assignment. (iii) \emph{What the rule missed.} T-wave abnormality (164934002; STTC share 0.88) is the most frequent near miss; abnormal QRS and Q-wave abnormality reach 0.80--0.81 for MI, but PTB-XL treats the corresponding statements as descriptions of form, not as diagnoses, so V2 and V4, which add all three, are deliberately liberal. Codes that never occur in the organizers' PTB-XL encoding cannot be assigned by a data-driven rule at all: right bundle branch block (59118001) and complete left bundle branch block (733534002), which the organizers score as equivalent to assigned codes (variant V5 adds them), and left ventricular high voltage (55827005; China only), which stays unmapped in every variant because assigning it to HYP would be a clinical judgement that we are not qualified to make.

\begingroup\scriptsize
\setlength{\tabcolsep}{3pt}
\begin{longtable}{lp{4.8cm}llrrrr}
\caption{Additional analysis: SNOMED CT codes of included and excluded records, Georgia, USA.}\label{tab:s_codes_georgia}\\
\toprule
SNOMED CT & Clinical name (abbreviation) & Assigned & Largest share & Included & \% & Excluded & \% \\
\midrule
\endfirsthead
\toprule
SNOMED CT & Clinical name (abbreviation) & Assigned & Largest share & Included & \% & Excluded & \% \\
\midrule
\endhead
\bottomrule
\endfoot
164934002 & t wave abnormal (TAb) & -- & 0.88 STTC & 1{,}857 & 30.7 & 442 & 10.4 \\
428750005 & nonspecific st t abnormality (NSSTTA) & STTC & 1.00 STTC & 1{,}874 & 31.0 & 0 & 0.0 \\
426783006 & sinus rhythm (NSR) & -- & 0.22 MI & 0 & 0.0 & 1{,}735 & 40.9 \\
426177001 & sinus bradycardia (SB) & -- & 0.21 CD & 906 & 15.0 & 762 & 18.0 \\
111975006 & prolonged qt interval (LQT) & STTC & 0.99 STTC & 1{,}383 & 22.9 & 0 & 0.0 \\
427084000 & sinus tachycardia (STach) & -- & 0.36 MI & 751 & 12.4 & 510 & 12.0 \\
164873001 & left ventricular hypertrophy (LVH) & HYP & 0.91 HYP & 1{,}225 & 20.3 & 0 & 0.0 \\
164930006 & st interval abnormal (STIAb) & -- & not in PTB-XL & 967 & 16.0 & 18 & 0.4 \\
39732003 & left axis deviation (LAD) & -- & 0.45 CD & 657 & 10.9 & 281 & 6.6 \\
425623009 & lateral ischaemia (LIs) & STTC & 1.00 STTC & 901 & 14.9 & 0 & 0.0 \\
67741000119109 & left atrial enlargement (LAE) & HYP & 1.00 HYP & 862 & 14.3 & 0 & 0.0 \\
59931005 & t wave inversion (TInv) & -- & 0.66 STTC & 657 & 10.9 & 151 & 3.6 \\
270492004 & 1st degree av block (IAVB) & CD & 1.00 CD & 766 & 12.7 & 0 & 0.0 \\
284470004 & premature atrial contraction (PAC) & -- & 0.35 STTC & 406 & 6.7 & 230 & 5.4 \\
164889003 & atrial fibrillation (AF) & -- & 0.60 STTC & 364 & 6.0 & 204 & 4.8 \\
59118001 & right bundle branch block (RBBB) & -- & not in PTB-XL & 312 & 5.2 & 228 & 5.4 \\
164917005 & qwave abnormal (QAb) & -- & 0.80 MI & 416 & 6.9 & 47 & 1.1 \\
427393009 & sinus arrhythmia (SA) & -- & 0.21 MI & 224 & 3.7 & 230 & 5.4 \\
425419005 & inferior ischaemia (IIs) & STTC & 1.00 STTC & 449 & 7.4 & 0 & 0.0 \\
713426002 & incomplete right bundle branch block (IRBBB) & CD & 1.00 CD & 402 & 6.6 & 0 & 0.0 \\
251146004 & low qrs voltages (LQRSV) & -- & 0.27 STTC & 169 & 2.8 & 204 & 4.8 \\
17338001 & ventricular premature beats (VPB) & -- & not in PTB-XL & 257 & 4.3 & 99 & 2.3 \\
\end{longtable}
\endgroup
{\footnotesize The 22 most frequent of 67 distinct codes in this cohort (42 of them occur in the organizers' PTB-XL encoding and could therefore be mapped). Included: 6{,}047 records; excluded for having no mapped superclass: 4{,}245 (41.0\% of all records). A record can carry several codes, so percentages do not sum to 100. Largest share: the largest share of a non-NORM superclass among PTB-XL records with that code. Of the excluded records, 2{,}586 carry only sinus rhythm or sinus-rate statements (sinus rhythm, bradycardia, tachycardia, arrhythmia; 1{,}735 only sinus rhythm), 447 carry a code with a share between 0.80 and 0.90, and 174 carry no code that occurs in PTB-XL. Mean codes per record: 2.95 (included), 1.31 (excluded).\par}
\medskip

\begingroup\scriptsize
\setlength{\tabcolsep}{3pt}
\begin{longtable}{lp{4.8cm}llrrrr}
\caption{Additional analysis: SNOMED CT codes of included and excluded records, China.}\label{tab:s_codes_china}\\
\toprule
SNOMED CT & Clinical name (abbreviation) & Assigned & Largest share & Included & \% & Excluded & \% \\
\midrule
\endfirsthead
\toprule
SNOMED CT & Clinical name (abbreviation) & Assigned & Largest share & Included & \% & Excluded & \% \\
\midrule
\endhead
\bottomrule
\endfoot
426177001 & sinus bradycardia (SB) & -- & 0.21 CD & 2{,}081 & 23.1 & 14{,}478 & 40.1 \\
426783006 & sinus rhythm (NSR) & -- & 0.22 MI & 693 & 7.7 & 7{,}432 & 20.6 \\
164890007 & atrial flutter (AFL) & -- & 0.21 STTC & 3{,}374 & 37.5 & 4{,}684 & 13.0 \\
427084000 & sinus tachycardia (STach) & -- & 0.36 MI & 1{,}792 & 19.9 & 5{,}463 & 15.1 \\
164934002 & t wave abnormal (TAb) & -- & 0.88 STTC & 1{,}427 & 15.8 & 5{,}616 & 15.5 \\
55827005 & left ventricular high voltage (LVHV) & -- & not in PTB-XL & 1{,}625 & 18.0 & 3{,}775 & 10.4 \\
55930002 & s t changes (STC) & STTC & 1.00 STTC & 4{,}231 & 47.0 & 0 & 0.0 \\
59931005 & t wave inversion (TInv) & -- & 0.66 STTC & 1{,}433 & 15.9 & 1{,}444 & 4.0 \\
427393009 & sinus arrhythmia (SA) & -- & 0.21 MI & 244 & 2.7 & 2{,}306 & 6.4 \\
164889003 & atrial fibrillation (AF) & -- & 0.60 STTC & 608 & 6.8 & 1{,}172 & 3.2 \\
429622005 & st depression (STD) & -- & 0.74 STTC & 1{,}448 & 16.1 & 220 & 0.6 \\
39732003 & left axis deviation (LAD) & -- & 0.45 CD & 473 & 5.3 & 1{,}071 & 3.0 \\
284470004 & premature atrial contraction (PAC) & -- & 0.35 STTC & 510 & 5.7 & 802 & 2.2 \\
10370003 & pacing rhythm (PR) & -- & 0.01 MI & 258 & 2.9 & 923 & 2.6 \\
428750005 & nonspecific st t abnormality (NSSTTA) & STTC & 1.00 STTC & 1{,}158 & 12.9 & 0 & 0.0 \\
270492004 & 1st degree av block (IAVB) & CD & 1.00 CD & 1{,}140 & 12.7 & 0 & 0.0 \\
713427006 & complete right bundle branch block (CRBBB) & CD & 1.00 CD & 1{,}095 & 12.2 & 0 & 0.0 \\
427172004 & premature ventricular contractions (PVC) & -- & not in PTB-XL & 440 & 4.9 & 651 & 1.8 \\
164917005 & qwave abnormal (QAb) & -- & 0.80 MI & 435 & 4.8 & 628 & 1.7 \\
251146004 & low qrs voltages (LQRSV) & -- & 0.27 STTC & 256 & 2.8 & 787 & 2.2 \\
47665007 & right axis deviation (RAD) & -- & 0.59 CD & 243 & 2.7 & 610 & 1.7 \\
164930006 & st interval abnormal (STIAb) & -- & not in PTB-XL & 242 & 2.7 & 559 & 1.5 \\
\end{longtable}
\endgroup
{\footnotesize The 22 most frequent of 94 distinct codes in this cohort (41 of them occur in the organizers' PTB-XL encoding and could therefore be mapped). Included: 9{,}004 records; excluded for having no mapped superclass: 36{,}146 (80.1\% of all records). A record can carry several codes, so percentages do not sum to 100. Largest share: the largest share of a non-NORM superclass among PTB-XL records with that code. Of the excluded records, 19{,}851 carry only sinus rhythm or sinus-rate statements (sinus rhythm, bradycardia, tachycardia, arrhythmia; 5{,}908 only sinus rhythm), 5{,}933 carry a code with a share between 0.80 and 0.90, and 433 carry no code that occurs in PTB-XL. Mean codes per record: 3.40 (included), 1.59 (excluded).\par}
\medskip

\begingroup\scriptsize
\setlength{\tabcolsep}{3pt}
\begin{longtable}{llrrrrr}
\caption{Additional analysis: records and positives (prevalence, \%) per superclass under each label-mapping variant.}\label{tab:s_prevalence}\\
\toprule
Cohort & Variant & Records & MI & STTC & CD & HYP \\
\midrule
\endfirsthead
\toprule
Cohort & Variant & Records & MI & STTC & CD & HYP \\
\midrule
\endhead
\bottomrule
\endfoot
PTB-XL test & - & 2{,}158 & 550 (25.5) & 521 (24.1) & 496 (23.0) & 262 (12.1) \\
\midrule
Georgia, USA & V0 & 6{,}047 & 7 (0.1)$^*$ & 4{,}059 (67.1) & 1{,}823 (30.1) & 2{,}088 (34.5) \\
Georgia, USA & V1 & 5{,}489 & 7 (0.1)$^*$ & 4{,}059 (73.9) & 1{,}769 (32.2) & 961 (17.5) \\
Georgia, USA & V2 & 6{,}047 & 423 (7.0) & 4{,}238 (70.1) & 1{,}823 (30.1) & 2{,}088 (34.5) \\
Georgia, USA & V3 & 10{,}292 & 7 (0.1)$^*$ & 4{,}059 (39.4) & 1{,}823 (17.7) & 2{,}088 (20.3) \\
Georgia, USA & V4 & 10{,}292 & 470 (4.6) & 4{,}680 (45.5) & 1{,}823 (17.7) & 2{,}088 (20.3) \\
Georgia, USA & V5 & 10{,}292 & 7 (0.1)$^*$ & 4{,}059 (39.4) & 2{,}203 (21.4) & 2{,}088 (20.3) \\
\midrule
China & V0 & 9{,}004 & 169 (1.9) & 5{,}670 (63.0) & 3{,}857 (42.8) & 770 (8.6) \\
China & V1 & 8{,}875 & 169 (1.9) & 5{,}670 (63.9) & 3{,}790 (42.7) & 131 (1.5) \\
China & V2 & 9{,}004 & 542 (6.0) & 6{,}574 (73.0) & 3{,}857 (42.8) & 770 (8.6) \\
\end{longtable}
\endgroup
{\footnotesize V0: published cohort and 90\% rule. V1: share $\geq$0.99 (records left without a label dropped). V2: share $\geq$0.80 on the published cohort. V3: published cohort plus every record excluded only for having no mapped superclass, counted as negative for all labels. V4: as V3 with share $\geq$0.80. V5: as V3, and the codes that the organizers score as equivalent to an assigned code (right bundle branch block 59118001 = complete right bundle branch block; complete left bundle branch block 733534002 = left bundle branch block) inherit its assignment. $^*$Not eligible in that variant (fewer than 50 positives or 50 negatives). NORM has no code in any variant.\par}
\medskip

\begingroup\scriptsize
\setlength{\tabcolsep}{3pt}
\begin{longtable}{llrrrrrr}
\caption{Additional analysis: external macro AUROC (mean over seeds) under each label-mapping variant, Georgia, USA.}\label{tab:s_mapping_perf_georgia}\\
\toprule
Pages & Model & V0 & V1 & V2 & V3 & V4 & V5 \\
\midrule
\endfirsthead
\toprule
Pages & Model & V0 & V1 & V2 & V3 & V4 & V5 \\
\midrule
\endhead
\bottomrule
\endfoot
Adaptive & CNN-only & 0.652 & 0.631 & 0.662 & 0.711 & 0.724 & 0.729 \\
Adaptive & CNN-only, capacity-matched & 0.700 & 0.672 & 0.710 & 0.746 & 0.762 & 0.765 \\
Adaptive & CNN-LSTM & 0.723 & 0.690 & 0.734 & 0.766 & 0.781 & 0.785 \\
Adaptive & CNN + temporal conv. & 0.707 & 0.685 & 0.718 & 0.759 & 0.775 & 0.777 \\
Adaptive & CNN + small ViT & 0.730 & 0.696 & 0.740 & 0.769 & 0.785 & 0.788 \\
Adaptive & ResNet-18, ImageNet & 0.773 & 0.738 & 0.784 & 0.802 & 0.819 & 0.818 \\
\midrule
Fixed & CNN-only & 0.707 & 0.674 & 0.717 & 0.752 & 0.766 & 0.769 \\
Fixed & CNN-only, capacity-matched & 0.733 & 0.696 & 0.744 & 0.773 & 0.788 & 0.790 \\
Fixed & CNN-LSTM & 0.738 & 0.700 & 0.749 & 0.777 & 0.793 & 0.796 \\
Fixed & CNN + temporal conv. & 0.738 & 0.703 & 0.748 & 0.780 & 0.795 & 0.798 \\
Fixed & CNN + small ViT & 0.737 & 0.700 & 0.748 & 0.777 & 0.793 & 0.795 \\
Fixed & ResNet-18, ImageNet & 0.779 & 0.740 & 0.790 & 0.812 & 0.829 & 0.827 \\
\midrule
Adaptive & PCA + LDA & 0.667 & 0.633 & 0.673 & 0.701 & 0.708 & 0.717 \\
\midrule
Fixed & PCA + LDA & 0.660 & 0.620 & 0.664 & 0.699 & 0.702 & 0.714 \\
\midrule
Signal & 1D ResNet (raw signals) & 0.787 & 0.753 & 0.797 & 0.818 & 0.833 & 0.831 \\
\end{longtable}
\endgroup
{\footnotesize V0: published cohort and 90\% rule. V1: share $\geq$0.99 (records left without a label dropped). V2: share $\geq$0.80 on the published cohort. V3: published cohort plus every record excluded only for having no mapped superclass, counted as negative for all labels. V4: as V3 with share $\geq$0.80. V5: as V3, and the codes that the organizers score as equivalent to an assigned code (right bundle branch block 59118001 = complete right bundle branch block; complete left bundle branch block 733534002 = left bundle branch block) inherit its assignment. Macro AUROC over the labels of the published analysis that remain eligible in the variant (V0: STTC, CD, HYP; V1: STTC, CD, HYP; V2: STTC, CD, HYP; V3: STTC, CD, HYP; V4: STTC, CD, HYP; V5: STTC, CD, HYP). V0 reproduces Table~4 of the main text. Predictions of the published cohort are the saved ones; the added records received one inference per run.\par}
\medskip

\begingroup\scriptsize
\setlength{\tabcolsep}{3pt}
\begin{longtable}{llrrr}
\caption{Additional analysis: external macro AUROC (mean over seeds) under each label-mapping variant, China.}\label{tab:s_mapping_perf_china}\\
\toprule
Pages & Model & V0 & V1 & V2 \\
\midrule
\endfirsthead
\toprule
Pages & Model & V0 & V1 & V2 \\
\midrule
\endhead
\bottomrule
\endfoot
Adaptive & CNN-only & 0.717 & 0.642 & 0.742 \\
Adaptive & CNN-only, capacity-matched & 0.777 & 0.687 & 0.794 \\
Adaptive & CNN-LSTM & 0.799 & 0.708 & 0.812 \\
Adaptive & CNN + temporal conv. & 0.811 & 0.720 & 0.823 \\
Adaptive & CNN + small ViT & 0.815 & 0.716 & 0.824 \\
Adaptive & ResNet-18, ImageNet & 0.848 & 0.781 & 0.865 \\
\midrule
Fixed & CNN-only & 0.768 & 0.669 & 0.787 \\
Fixed & CNN-only, capacity-matched & 0.809 & 0.711 & 0.821 \\
Fixed & CNN-LSTM & 0.817 & 0.725 & 0.830 \\
Fixed & CNN + temporal conv. & 0.824 & 0.723 & 0.835 \\
Fixed & CNN + small ViT & 0.819 & 0.716 & 0.831 \\
Fixed & ResNet-18, ImageNet & 0.853 & 0.775 & 0.867 \\
\midrule
Adaptive & PCA + LDA & 0.732 & 0.652 & 0.740 \\
\midrule
Fixed & PCA + LDA & 0.752 & 0.644 & 0.743 \\
\midrule
Signal & 1D ResNet (raw signals) & 0.855 & 0.791 & 0.873 \\
\end{longtable}
\endgroup
{\footnotesize V0: published cohort and 90\% rule. V1: share $\geq$0.99 (records left without a label dropped). V2: share $\geq$0.80 on the published cohort. V3: published cohort plus every record excluded only for having no mapped superclass, counted as negative for all labels. V4: as V3 with share $\geq$0.80. V5: as V3, and the codes that the organizers score as equivalent to an assigned code (right bundle branch block 59118001 = complete right bundle branch block; complete left bundle branch block 733534002 = left bundle branch block) inherit its assignment. Macro AUROC over the labels of the published analysis that remain eligible in the variant (V0: MI, STTC, CD, HYP; V1: MI, STTC, CD, HYP; V2: MI, STTC, CD, HYP). V0 reproduces Table~4 of the main text. Predictions of the published cohort are the saved ones; the added records received one inference per run.\par}
\medskip

\begingroup\scriptsize
\setlength{\tabcolsep}{3pt}
\begin{longtable}{llrrrccrr}
\caption{Additional analysis: the external hypotheses H1 and H2 under each label-mapping variant, Georgia, USA: difference in mean-over-seeds AUROC.}\label{tab:s_mapping_hyp_georgia}\\
\toprule
Family & Comparison & V0 & V1 & V2 & V3 [95\% CI] & V3 $p$ & V4 & V5 \\
\midrule
\endfirsthead
\toprule
Family & Comparison & V0 & V1 & V2 & V3 [95\% CI] & V3 $p$ & V4 & V5 \\
\midrule
\endhead
\bottomrule
\endfoot
H1 & Capacity-m. $-$ CNN-only & +0.048 & +0.040 & +0.048 & +0.035 [+0.023, +0.047] & $\le$0.002 & +0.038 & +0.037 \\
H1 & CNN-LSTM $-$ CNN-only & +0.071 & +0.059 & +0.072 & +0.055 [+0.047, +0.064] & $\le$0.002 & +0.057 & +0.056 \\
H1 & Temporal conv. $-$ CNN-only & +0.055 & +0.053 & +0.056 & +0.048 [+0.042, +0.054] & $\le$0.002 & +0.051 & +0.048 \\
H1 & ViT $-$ CNN-only & +0.078 & +0.065 & +0.079 & +0.059 [+0.051, +0.065] & $\le$0.002 & +0.061 & +0.059 \\
H1 & CNN-LSTM $-$ Capacity-m. & +0.023 & +0.019 & +0.024 & +0.020 [+0.009, +0.030] & $\le$0.002 & +0.020 & +0.019 \\
H1 & Temporal conv. $-$ Capacity-m. & +0.007 & +0.013 & +0.008 & +0.013 [+0.001, +0.025] & 0.070 & +0.014 & +0.011 \\
H1 & ViT $-$ Capacity-m. & +0.030 & +0.024 & +0.030 & +0.023 [+0.012, +0.035] & $\le$0.002 & +0.024 & +0.022 \\
H1 & Temporal conv. $-$ CNN-LSTM & $-$0.016 & $-$0.005 & $-$0.015 & $-$0.007 [$-$0.013, $-$0.001] & 0.070 & $-$0.006 & $-$0.008 \\
H1 & ViT $-$ CNN-LSTM & +0.007 & +0.006 & +0.007 & +0.004 [$-$0.006, +0.015] & 0.524 & +0.004 & +0.003 \\
H1 & ViT $-$ Temporal conv. & +0.023 & +0.011 & +0.022 & +0.010 [+0.002, +0.019] & 0.070 & +0.010 & +0.011 \\
\midrule
H2 & CNN-only, macro & +0.055 & +0.043 & +0.056 & +0.041 [+0.034, +0.048] & $\le$0.002 & +0.042 & +0.041 \\
H2 & CNN-only, HYP & +0.129 & +0.079 & +0.129 & +0.094 [+0.081, +0.108] & $\le$0.002 & +0.094 & +0.094 \\
H2 & CNN-LSTM, macro & +0.015 & +0.010 & +0.015 & +0.011 [+0.002, +0.020] & 0.015 & +0.012 & +0.011 \\
H2 & CNN-LSTM, HYP & +0.027 & +0.016 & +0.027 & +0.022 [+0.007, +0.038] & 0.004 & +0.022 & +0.022 \\
H2 & Temporal conv., macro & +0.031 & +0.019 & +0.030 & +0.021 [+0.014, +0.028] & $\le$0.002 & +0.020 & +0.021 \\
H2 & Temporal conv., HYP & +0.096 & +0.057 & +0.096 & +0.067 [+0.054, +0.082] & $\le$0.002 & +0.067 & +0.067 \\
H2 & ViT, macro & +0.008 & +0.004 & +0.007 & +0.007 [+0.004, +0.011] & 0.002 & +0.007 & +0.008 \\
H2 & ViT, HYP & +0.018 & +0.011 & +0.018 & +0.016 [+0.004, +0.029] & 0.007 & +0.016 & +0.016 \\
\end{longtable}
\endgroup
{\footnotesize V0: published cohort and 90\% rule. V1: share $\geq$0.99 (records left without a label dropped). V2: share $\geq$0.80 on the published cohort. V3: published cohort plus every record excluded only for having no mapped superclass, counted as negative for all labels. V4: as V3 with share $\geq$0.80. V5: as V3, and the codes that the organizers score as equivalent to an assigned code (right bundle branch block 59118001 = complete right bundle branch block; complete left bundle branch block 733534002 = left bundle branch block) inherit its assignment. H1: head pairs on adaptive-gain pages (macro AUROC). H2: fixed-gain minus adaptive-gain training (macro and HYP AUROC). V0 repeats the published point estimates. V3: hierarchical seed $\times$ record bootstrap (10{,}000 resamples), Holm within H1 (10) and H2 (8); the other variants are point estimates. $p$ printed as ``$\le$'' is at the resolution limit of the bootstrap: no replicate crossed zero, and no smaller value is claimed (Section~\ref{sec:s_review}). \par}
\medskip

\begingroup\scriptsize
\setlength{\tabcolsep}{3pt}
\begin{longtable}{llrrr}
\caption{Additional analysis: the external hypotheses H1 and H2 under each label-mapping variant, China: difference in mean-over-seeds AUROC.}\label{tab:s_mapping_hyp_china}\\
\toprule
Family & Comparison & V0 & V1 & V2 \\
\midrule
\endfirsthead
\toprule
Family & Comparison & V0 & V1 & V2 \\
\midrule
\endhead
\bottomrule
\endfoot
H1 & Capacity-m. $-$ CNN-only & +0.060 & +0.046 & +0.053 \\
H1 & CNN-LSTM $-$ CNN-only & +0.083 & +0.066 & +0.071 \\
H1 & Temporal conv. $-$ CNN-only & +0.094 & +0.078 & +0.082 \\
H1 & ViT $-$ CNN-only & +0.099 & +0.074 & +0.083 \\
H1 & CNN-LSTM $-$ Capacity-m. & +0.023 & +0.020 & +0.018 \\
H1 & Temporal conv. $-$ Capacity-m. & +0.034 & +0.032 & +0.029 \\
H1 & ViT $-$ Capacity-m. & +0.039 & +0.029 & +0.030 \\
H1 & Temporal conv. $-$ CNN-LSTM & +0.011 & +0.012 & +0.011 \\
H1 & ViT $-$ CNN-LSTM & +0.016 & +0.008 & +0.012 \\
H1 & ViT $-$ Temporal conv. & +0.005 & $-$0.004 & +0.001 \\
\midrule
H2 & CNN-only, macro & +0.051 & +0.027 & +0.045 \\
H2 & CNN-only, HYP & +0.112 & +0.014 & +0.112 \\
H2 & CNN-LSTM, macro & +0.018 & +0.017 & +0.017 \\
H2 & CNN-LSTM, HYP & +0.035 & +0.035 & +0.035 \\
H2 & Temporal conv., macro & +0.014 & +0.004 & +0.012 \\
H2 & Temporal conv., HYP & +0.033 & $-$0.005 & +0.033 \\
H2 & ViT, macro & +0.004 & $-$0.000 & +0.007 \\
H2 & ViT, HYP & +0.014 & +0.001 & +0.014 \\
\end{longtable}
\endgroup
{\footnotesize V0: published cohort and 90\% rule. V1: share $\geq$0.99 (records left without a label dropped). V2: share $\geq$0.80 on the published cohort. V3: published cohort plus every record excluded only for having no mapped superclass, counted as negative for all labels. V4: as V3 with share $\geq$0.80. V5: as V3, and the codes that the organizers score as equivalent to an assigned code (right bundle branch block 59118001 = complete right bundle branch block; complete left bundle branch block 733534002 = left bundle branch block) inherit its assignment. H1: head pairs on adaptive-gain pages (macro AUROC). H2: fixed-gain minus adaptive-gain training (macro and HYP AUROC). V0 repeats the published point estimates. V3: hierarchical seed $\times$ record bootstrap (10{,}000 resamples), Holm within H1 (10) and H2 (8); the other variants are point estimates. $p$ printed as ``$\le$'' is at the resolution limit of the bootstrap: no replicate crossed zero, and no smaller value is claimed (Section~\ref{sec:s_review}). \par}
\medskip

\subsection{Bootstrap $p$-values and Holm correction}

\begingroup\scriptsize
\setlength{\tabcolsep}{3pt}
\begin{longtable}{p{4.2cm}p{4.4cm}rrrlrr}
\caption{Additional analysis: bootstrap $p$-values before and after the statistical revision of 15 November 2026, by comparison family.}\label{tab:s_stats_families}\\
\toprule
Analysis & Family & $m$ & Old $p_\mathrm{adj}=0$ & At floor & Printed as & Old $<$0.05 & Revised $<$0.05 \\
\midrule
\endfirsthead
\toprule
Analysis & Family & $m$ & Old $p_\mathrm{adj}=0$ & At floor & Printed as & Old $<$0.05 & Revised $<$0.05 \\
\midrule
\endhead
\bottomrule
\endfoot
PTB-XL head pairs (S7) & test | 42-46 (primary) & 10 & 5 & 5 & $\le$0.002 & 7 & 8 \\
PTB-XL head pairs (S7) & test | 43-46 GPU only & 10 & 5 & 5 & $\le$0.002 & 7 & 7 \\
PTB-XL head pairs (S7) & validation | 42-46 (primary) & 10 & 6 & 6 & $\le$0.002 & 8 & 8 \\
PTB-XL head pairs (S7) & validation | 43-46 GPU only & 10 & 6 & 6 & $\le$0.002 & 7 & 7 \\
PTB-XL fixed $-$ adaptive (Table 3, S8) & 42-46 & 8 & 4 & 4 & $\le$0.002 & 5 & 5 \\
PTB-XL fixed $-$ adaptive (Table 3, S8) & 43-46 GPU only & 8 & 4 & 4 & $\le$0.002 & 4 & 4 \\
Georgia H1, H2 (Tables 2, 3) & H1 & 10 & 8 & 8 & $\le$0.002 & 8 & 8 \\
Georgia H1, H2 (Tables 2, 3) & H2 & 8 & 7 & 4 & $\le$0.002 & 8 & 8 \\
China H1, H2 (Tables 2, 3) & H1 & 10 & 8 & 8 & $\le$0.002 & 8 & 8 \\
China H1, H2 (Tables 2, 3) & H2 & 8 & 4 & 4 & $\le$0.002 & 7 & 7 \\
ResNet-18 fixed $-$ adaptive (S15) & PTB-XL test & 2 & 1 & 0 & $\le$0.001 & 2 & 2 \\
ResNet-18 fixed $-$ adaptive (S15) & external cohorts & 4 & 0 & 0 & $\le$0.001 & 0 & 0 \\
40-epoch sensitivity (S17) & 40 - 20 epochs & 5 & 4 & 4 & $\le$0.001 & 4 & 4 \\
40-epoch sensitivity (S17) & head pairs, 40 epochs & 10 & 4 & 4 & $\le$0.002 & 5 & 5 \\
40-epoch sensitivity (S17) & head pairs, 20 epochs (same seeds) & 10 & 6 & 6 & $\le$0.002 & 8 & 8 \\
Fixed gain, 40 epochs (S18) & fixed - adaptive, 40 epochs & 8 & 3 & 2 & $\le$0.002 & 3 & 3 \\
Fixed gain, 40 epochs (S18) & fixed gain, 40 - 20 epochs & 4 & 2 & 1 & $\le$0.001 & 2 & 2 \\
Fixed gain, 40 epochs (S18) & fixed - adaptive, 20 epochs (same seeds) & 8 & 4 & 4 & $\le$0.002 & 4 & 4 \\
\end{longtable}
\endgroup
{\footnotesize Old: 2{,}000 resamples, $p=2\min\{P(\Delta^*\le0),P(\Delta^*\ge0)\}$, which is 0 when no replicate crosses zero and stays 0 after Holm correction; such values were printed as ``$<$0.001''. Revised: 10{,}000 resamples and the add-one estimator $p=2\min\{(1+\#[\Delta^*\le0])/(B+1),(1+\#[\Delta^*\ge0])/(B+1)\}$, whose smallest value is 0.0002; after Holm correction in a family of $m$ comparisons the smallest attainable value is $m$ times that (column ``Printed as''). $m$: comparisons in the Holm family. Old $p_\mathrm{adj}=0$: comparisons whose old adjusted $p$ was exactly zero. At floor: comparisons with no replicate across zero in the revised bootstrap. Last two columns: comparisons with adjusted $p<0.05$. In total 143 comparisons; 81 had an old adjusted $p$ of zero; 1 changed side of 0.05. The first 2{,}000 replicates of the revised bootstrap are the old random draws and reproduce every archived old value exactly (checked). The old outputs are archived in results/stats\_legacy\_2000resamples\_20261115/; every comparison is listed in results/review\_20261115/stats\_revision/old\_vs\_revised.csv.\par}
\medskip

\begingroup\scriptsize
\setlength{\tabcolsep}{3pt}
\begin{longtable}{p{3.6cm}p{3.6cm}p{4.2cm}rrr}
\caption{Additional analysis: comparisons whose Holm-adjusted $p$ moved across or near 0.05 with the revised bootstrap.}\label{tab:s_stats_changed}\\
\toprule
Analysis & Family & Comparison & Difference & Old $p$ & Revised $p$ \\
\midrule
\endfirsthead
\toprule
Analysis & Family & Comparison & Difference & Old $p$ & Revised $p$ \\
\midrule
\endhead
\bottomrule
\endfoot
\multicolumn{6}{l}{\emph{Changed side of 0.05}} \\
PTB-XL head pairs (S7) & test | 42-46 (primary) & CNN + small ViT minus CNN-LSTM & +0.009 & 0.054 & 0.044 \\
\multicolumn{6}{l}{\emph{Other changes of at least 0.005 near 0.05}} \\
PTB-XL head pairs (S7) & test | 43-46 GPU only & CNN + temporal conv minus CNN-only, capacity-matched & +0.014 & 0.015 & 0.009 \\
PTB-XL head pairs (S7) & validation | 43-46 GPU only & CNN + temporal conv minus CNN-only, capacity-matched & +0.013 & 0.028 & 0.022 \\
PTB-XL fixed $-$ adaptive (Table 3, S8) & 42-46 & cnn\_lstm | macro AUROC & +0.012 & 0.028 & 0.017 \\
PTB-XL fixed $-$ adaptive (Table 3, S8) & 43-46 GPU only & cnn\_lstm | macro AUROC & +0.012 & 0.100 & 0.084 \\
ResNet-18 fixed $-$ adaptive (S15) & external cohorts & georgia | HYP AUROC & +0.011 & 0.084 & 0.076 \\
40-epoch sensitivity (S17) & head pairs, 20 epochs (same seeds) & cnn\_lstm - cnn\_pool\_wide | test macro AUROC & +0.014 & 0.042 & 0.024 \\
\end{longtable}
\endgroup
{\footnotesize Changes away from the resolution floor come from Monte Carlo error only (2{,}000 versus 10{,}000 resamples) and from the add-one term.\par}
\medskip

\section{Model cards}

All models: PTB-XL v1.0.3, official patient folds (1-8 train, 9 validation, 10 test), five overlapping superclasses (NORM, MI, STTC, CD, HYP), training-only normalization, AdamW lr 1e-3, batch 16, max 20 epochs, patience 4, cosine schedule, unweighted BCE, no augmentation, no pretrained weights (except the ResNet-18 sensitivity model), checkpoint selected on validation, threshold 0.5, one test inference. Fold 10 was observed before the models were specified, so test results are exploratory. Intended use: research on ECG-image modelling only, not clinical use. Known limitations: clean synthetic renderings (no scans/photos), one population (Germany), label noise, and adaptive page gain on v1 pages leaks amplitude (and therefore hypertrophy) information.

\begin{itemize}
\item \textbf{PCA + LDA / PCA + linear SVM.} Classical baselines on PCA of the downsampled page; one fit each (no seeds). Hyperparameters chosen on validation by a frozen grid.
\item \textbf{CNN-only (cnn\_pool).} Shared 1,220-parameter lead-image CNN, mean over 70 feature columns, linear readout. 1,705 parameters. Sees only local image regions; cannot relate the calibration pulse to the waveform.
\item \textbf{Wide CNN-only (cnn\_pool\_wide).} Same CNN, per-column MLP 96-96-64 then column mean. 17,065 parameters: capacity of the sequence heads without column mixing (protocol amendment 3 Nov 2026).
\item \textbf{CNN-LSTM (cnn\_lstm).} Same CNN, one unidirectional LSTM layer (hidden width 32) over the 70 columns, mean of outputs. 18,025 parameters.
\item \textbf{CNN + temporal convolution (temporal\_conv).} Same CNN, two causal dilated residual Conv1d blocks (receptive field 13 columns). 17,161 parameters.
\item \textbf{CNN + small ViT (hybrid\_vit).} Same CNN, small transformer encoder over the columns with positional embeddings. 23,881 parameters. Sees the whole row at once.
\item \textbf{Signal ResNet1D (resnet1d\_wang).} Reference model on the raw 100 Hz waveforms instead of images. About 510k parameters. Shows what the image representation costs.
\item \textbf{ImageNet-pretrained ResNet-18 (resnet18).} Sensitivity model outside the controlled comparison (amendment 11 Nov 2026): torchvision's ImageNet ResNet-18 weights, the 12 lead rows stacked into one 1440x1120 image, first-layer RGB filters summed, global average pooling, linear readout. 11.2 million parameters, fine-tuned with the shared settings in bfloat16 mixed precision; seeds 42-44 on adaptive and fixed gain.
\end{itemize}
Every completed run, with hardware, selected epoch and results, is listed in Supplementary Table~\ref{tab:s_runs}. Generated from src/make\_model\_cards.py.

\section{Protocol amendments}

\label{sec:s_amendments}\subsection{Amendment of 3 November 2026: seeds 45/46 and a capacity-matched head}


Written and frozen \textbf{before} any run listed here started. Nothing below changes a completed run,
a hyperparameter or a config used by an existing run. Motivation: \nolinkurl{docs/ROADMAP_Q1.md} weaknesses \#3
(three seeds give a rough SD; many comparisons without correction) and \#4 (CNN-only has 1,705 parameters,
the sequence heads 17,161--23,881).
\paragraph{A. Additional seeds}
\begin{itemize}
\item Seeds \textbf{45 and 46} for every network, under the unchanged configs:
\begin{itemize}
\item adaptive gain: \texttt{ptbxl\_v2\_\{cnn\_pool, cnn\_lstm, temporal\_conv, hybrid\_vit\}.json};
\item fixed gain: \texttt{ptbxl\_v2\_fg\_\{cnn\_pool, cnn\_lstm, temporal\_conv, hybrid\_vit\}.json}.
\end{itemize}
\item All runs use the PC GPU (\texttt{--device cuda}), the unchanged shared trainer (\nolinkurl{src/training/ptbxl_trainer.py})
and \nolinkurl{src/run_ptbxl_experiment.py}. These are the same code and settings as seeds 43/44.
\item Seeds 42--46 are all reported. No seed is dropped or selected. A failed run is reported as failed and gets
no rerun under the same ID.
\item Runs execute in up to five parallel lanes. Runtimes are therefore not speed measurements.
\end{itemize}
\paragraph{B. Capacity-matched CNN-only head (\nolinkurl{cnn_pool_wide})}
\begin{itemize}
\item Model: \nolinkurl{src/models/ptbxl_cnn_pool_wide.py}, config \nolinkurl{experiments/configs/ptbxl_v2_cnn_pool_wide.json},
runner \nolinkurl{src/run_ptbxl_wide_head.py}.
\item Architecture: the shared 1,220-parameter CNN, then one MLP (96 $\rightarrow$ 96 $\rightarrow$ 64, ReLU, dropout 0.2) applied to
each of the 70 feature columns separately, then the column mean and a linear 64 $\rightarrow$ 5 readout.
\textbf{17,065 parameters} in total, against 17,161 for temporal conv and 18,025 for CNN-LSTM.
\item No operation mixes columns before the mean. The head therefore has the capacity of the sequence heads
without their ability to relate one part of a row to another.
\item Hyperparameters are identical to every other network: AdamW lr 1e-3, batch 16, 20 epochs, patience 4,
cosine schedule, dropout 0.2. There is no tuning.
\item Seeds 42, 43, 44, 45, 46, all on the PC GPU, with adaptive-gain (v1) pages only.
\end{itemize}
\paragraph{C. Analysis (fixed now)}
\begin{enumerate}
\item \textbf{Main table.} For each network, report the mean $\pm$ SD across seeds of test and validation macro AUROC.
\begin{itemize}
\item Primary: seeds 42--46, with Mac-CPU seed 42 pooled, as in \nolinkurl{seed_summary_all_seeds.csv}.
\item Sensitivity: GPU-only seeds 43--46. The wide head has GPU seeds 42--46 in both.
\end{itemize}
\item \textbf{Pairwise head comparisons.} Pairs are taken among CNN-only, wide CNN-only, CNN-LSTM, temporal conv and
hybrid ViT (10 pairs).
\begin{itemize}
\item \textbf{Method.} For each pair, run a hierarchical bootstrap (2,000 resamples, seed 2024): resample seeds with
replacement, then patients with replacement, both shared across the two models. Report the mean
difference in test macro AUROC with a percentile 95\% CI.
\item \textbf{p-values.} Two-sided bootstrap p = 2$\cdot$min(P($\Delta$ $\leq$ 0), P($\Delta$ $\geq$ 0)).
\item \textbf{Correction.} Holm correction across the 10 pairs at $\alpha$ = 0.05.
\item \textbf{Status.} Fold 10 was observed before, so these are labelled exploratory. Validation results are
reported alongside.
\end{itemize}
\item \textbf{Interpretation rule for the wide head.}
\begin{itemize}
\item Wide head > CNN-only: extra head capacity helps.
\item Temporal conv or CNN-LSTM > wide head: relating columns helps beyond capacity.
\item The ViT is compared too, but it has about 6,800 more parameters than the wide head.
\end{itemize}
\item \textbf{Fixed gain.} \nolinkurl{src.compare_fixed_gain} is extended to seeds 45/46. The matched-hardware summary uses GPU
seeds 43--46, with the same metrics and strata as before.
\item \textbf{Limitations stated up front.}
\begin{itemize}
\item GPU training is not bit-reproducible across runs.
\item Seed 42 of the four original networks was trained on a Mac CPU.
\item The wide head also differs in having one extra dropout layer.
\end{itemize}
\end{enumerate}

\subsection{External-test amendment of 6 November 2026}

Recorded in experiments/external/amendments.json; the hash chain is checked before every external step.

\begin{itemize}\item Time (UTC): 2026-11-06T16:25:28+00:00; file: docs/external\_test\_plan.md.\item Exclude 2 Ningbo records whose signal files never match PhysioNet published SHA-256 (3 download attempts). Recorded after the Georgia inventory and before any Chapman-Shaoxing/Ningbo header was parsed.\end{itemize}

\subsection{Amendment of 11 November 2026: signal-reference seeds 45 and 46}


Written \textbf{before} either run below started and before any of their results existed
(2026-11-10 20:55 UTC = 2026-11-11 00:55 local time). Nothing below changes a completed run, a hyperparameter,
a config used by an existing run, or the frozen external-test analysis.
\paragraph{Reason}

The image networks were trained with seeds 42--46 (\nolinkurl{docs/protocol_amendment_2026-11-03.md}), but the
signal-domain reference (\nolinkurl{experiments/configs/ptbxl_v2_signal_resnet1d.json}) only with seeds 42--44, because
that amendment did not cover it. A reviewer-style read of the manuscript (11 Nov) flagged the mismatch. The two
seeds are added only so that every trained network in the paper reports the same seeds.
\paragraph{A. Additional runs}
\begin{itemize}
\item \nolinkurl{ptbxl_v2_signal_resnet1d_seed45} and \nolinkurl{ptbxl_v2_signal_resnet1d_seed46}.
\item Same frozen config (\nolinkurl{experiments/configs/ptbxl_v2_signal_resnet1d.json}, unchanged; its \nolinkurl{seeds} field
records the original plan and is not edited), same runner (\nolinkurl{src/run_ptbxl_signal.py}, unchanged), same
unchanged shared trainer, same device (the PC GPU, \nolinkurl{cuda}, as for seeds 42--44), same command as seeds 42--44:
\texttt{python -X utf8 -m src.run\_ptbxl\_signal --seed <seed>}.
\item One test inference per run, as for every run. A failed run is reported as failed and is not rerun under the
same ID; an interrupted run is moved to \nolinkurl{experiments/abandoned_runs/} and restarted from scratch.
\end{itemize}
\paragraph{B. Reporting on PTB-XL}
\begin{itemize}
\item The signal reference is reported as mean $\pm$ SD over seeds 42--46 wherever the image models report seeds 42--46.
The per-seed values of all five seeds are reported in the Supplementary Information, so the earlier
three-seed summary (0.928 $\pm$ 0.001) stays recoverable.
\item The paired signal-versus-best-image comparison at seed 42 is unchanged (it uses seed 42 only).
\end{itemize}
\paragraph{C. External test (descriptive sensitivity analysis)}
\begin{itemize}
\item The frozen external plan (\nolinkurl{docs/external_test_plan.md}, \S{}5) includes "every completed run", which at the
freeze meant signal seeds 42--44. It is \textbf{not} edited.
\item Seeds 45/46 get exactly one inference per external dataset (Georgia; Chapman-Shaoxing + Ningbo), with the
same pages/signals cache, the same 100 Hz resampling and the same one-inference guard
(\nolinkurl{inference_state.json}). Before inference, each reloaded model must reproduce its saved PTB-XL validation
probabilities (tolerance 1e-3), as for all other runs.
\item \nolinkurl{src/external_test.py} refuses to start inference while any run already has a state file. A new step
(\nolinkurl{infer-added}) runs only completed runs that have no state file yet and refuses to touch any other run. No
existing prediction file is rewritten.
\item These two runs are a \textbf{descriptive} addition. The confirmatory hypotheses H1 and H2 involve only the image
heads and are unchanged; re-running \nolinkurl{analyze} must reproduce \nolinkurl{hypotheses.csv} exactly, which is checked.
The signal reference stays a secondary, descriptive comparison, as in the frozen plan.
\end{itemize}

\subsection{Amendment of 11 November 2026: 40-epoch sensitivity and a pretrained ResNet-18}


Written and frozen \textbf{before} any run listed here started (11 Nov 2026, about 01:00 local time, UTC+4 = 10 Nov 21:00 UTC; runs started 21:05 UTC).
Nothing below changes a completed run, or a config or hyperparameter used by an existing run. Both
analyses are \textbf{sensitivity analyses}. They do not replace the primary results reported so far.

\textbf{Motivation.} On 11 Nov, a reviewer-style read of the manuscript raised two points:
\begin{enumerate}
\item \textbf{Possible under-training.} 17 of the 25 adaptive-gain network runs (seeds 42--46) selected their checkpoint
at epoch 18--20 of the 20-epoch budget. CNN-only and the capacity-matched head did so in every run. A reviewer
could argue that the simpler heads were still improving, which would bias the head comparison. This was
observed from the selected epochs (validation-based) of existing runs. No new test result has been seen.
\item \textbf{Very small models.} All networks have at most 24k parameters. It is open whether the adaptive-gain cue,
and the benefit of fixed gain, also apply to a standard, ImageNet-pretrained image backbone.
\end{enumerate}
\paragraph{A. 40-epoch sensitivity (\nolinkurl{ptbxl_v2_long40_*})}
\begin{itemize}
\item \textbf{Heads:}
\begin{itemize}
\item CNN-only (\nolinkurl{cnn_pool});
\item capacity-matched CNN-only (\nolinkurl{cnn_pool_wide});
\item CNN-LSTM (\nolinkurl{cnn_lstm});
\item CNN + temporal convolution (\nolinkurl{temporal_conv});
\item CNN + small ViT (\nolinkurl{hybrid_vit}).
\end{itemize}
\item \textbf{Pages and seeds:} adaptive-gain pages only (v1 cache); seeds \textbf{42, 43, 44}; all on the PC GPU.
\item \textbf{Configs:} \nolinkurl{experiments/configs/ptbxl_v2_long40_<head>.json}. Each is a copy of the head's existing v2 config
with only these changes:
\begin{itemize}
\item \nolinkurl{max_epochs} 40 instead of 20;
\item the cosine schedule therefore spans 40 epochs (\texttt{T\_max = max\_epochs}, as in every run);
\item new protocol, result family (\nolinkurl{ptbxl_v2_long40}) and run IDs.
\end{itemize}
\item \textbf{Unchanged:} everything else, including patience 4, lr 1e-3, batch 16, the selection rule, threshold 0.5,
one test inference per run, and the shared trainer and runner.
\item \textbf{Analysis (fixed now):}
\begin{enumerate}
\item Selected epoch and epochs run per run; how many runs stop early.
\item For each head, test and validation macro AUROC (mean $\pm$ SD over seeds 42--44) at 40 epochs, next to the
20-epoch runs of the same seeds. Seed 42 of the four original 20-epoch heads was CPU-trained; a GPU-only
view (seeds 43--44) is also shown.
\item The 10 pairwise head comparisons repeated on the 40-epoch runs. Method: the same hierarchical bootstrap as
the main analysis (2,000 resamples, seed 2024, seeds then patients, Holm over 10 pairs), with seeds 42--44.
\item \textbf{Interpretation rule.}
\begin{itemize}
\item If every sequence head still exceeds the capacity-matched head (Holm-adjusted p < 0.05), the claim
``relating columns helps beyond capacity'' is robust to the training budget.
\item If a gap shrinks below significance or reverses, the manuscript reports this and weakens the claim
accordingly.
\item Either outcome is reported.
\end{itemize}
\end{enumerate}
\end{itemize}
\paragraph{B. Pretrained ResNet-18 on adaptive and fixed gain (\nolinkurl{ptbxl_v2_resnet18}, \nolinkurl{ptbxl_v2_fg_resnet18})}
\begin{itemize}
\item \textbf{Weights:} torchvision's official ImageNet ResNet-18, \nolinkurl{resnet18-f37072fd.pth}, downloaded on 11 Nov 2026 from
\nolinkurl{https://download.pytorch.org/models/} with the user's permission.
\begin{itemize}
\item SHA-256 \nolinkurl{f37072fd47e89c5e827621c5baffa7500819f7896bbacec160b1a16c560e07ec}, stored in \nolinkurl{models/pretrained/}
and not versioned.
\item The network is re-implemented in \nolinkurl{src/models/ptbxl_resnet18.py}, with the same layer names as torchvision
(torchvision is not installed), and the weights are loaded with strict key matching.
\end{itemize}
\item \textbf{Input:}
\begin{itemize}
\item The same 12 $\times$ 120 $\times$ 1120 standardized ink tensor as every other network (same caches, same training-only
normalization).
\item The 12 lead rows are stacked vertically into one 1440 $\times$ 1120 single-channel page image. No resizing.
\item The first convolution's ImageNet RGB filters are summed into one input channel, which is equivalent to giving
the same image to all three channels.
\end{itemize}
\item \textbf{Head:} global average pooling, dropout 0.2, linear 512 $\rightarrow$ 5. \textbf{11.2 million parameters.}
\item \textbf{Training:}
\begin{itemize}
\item Identical to every other network: AdamW lr 1e-3, weight decay 1e-4, batch 16, max 20 epochs, patience 4,
cosine to 1e-5, gradient clipping 1.0, unweighted BCE, no augmentation, TF32 off, deterministic algorithms
requested.
\item \textbf{One deviation: bfloat16 mixed precision.} It applies to the forward pass on the GPU (autocast inside the
model) and uses a channels-last memory layout.
\begin{itemize}
\item Weights, optimizer state and loss stay float32, and logits are returned as float32.
\item Reason: in full float32, one epoch on 1440 $\times$ 1120 images took about 19 minutes (measured before any
training: 1.05 s per batch). With bf16 it takes about 3 minutes (0.17 s per batch). Only with bf16 do the
six runs fit in the time the shared PC is available.
\item This was decided before any ResNet run. The small networks keep full float32.
\end{itemize}
\item All weights are fine-tuned, and BatchNorm statistics update in training mode.
\item No hyperparameter is tuned.
\end{itemize}
\item \textbf{Seeds:} 42, 43, 44, on adaptive-gain and fixed-gain (2.5 mm/mV) pages, all on the PC GPU.
\item \textbf{Analysis (fixed now):}
\begin{enumerate}
\item PTB-XL test macro and per-label AUROC (mean $\pm$ SD over seeds), next to the small networks.
\item \textbf{Fixed $-$ adaptive} macro AUROC and HYP AUROC on the PTB-XL test fold. Method: the same hierarchical seed $\times$
patient bootstrap as the main fixed-gain analysis (seeds 42--44), Holm over these 2 comparisons.
\item \textbf{External cohorts (descriptive extension of the frozen external plan).}
\begin{itemize}
\item Each ResNet run is reloaded and must reproduce its saved validation probabilities.
\item It then performs exactly one inference per cohort (Georgia; Chapman-Shaoxing + Ningbo), with the frozen
inclusion rules, label mapping and eligibility.
\item Fixed $-$ adaptive macro and HYP AUROC: hierarchical seed $\times$ record bootstrap, Holm over the 4 comparisons
(2 cohorts $\times$ 2 metrics).
\item These are added after the plan was frozen, so they are reported as a prespecified-today sensitivity
analysis, not as part of the confirmatory H1/H2.
\end{itemize}
\item \textbf{Interpretation.}
\begin{itemize}
\item If fixed gain also improves HYP AUROC for ResNet-18, the gain cue is not specific to tiny models.
\item If not, the manuscript says so and restricts the claim.
\end{itemize}
\end{enumerate}
\end{itemize}
\paragraph{C. Execution and limits}
\begin{itemize}
\item \textbf{Order:} ResNet-18 runs first. They use the GPU heavily, and the PC is shared: the GPU must be free by
11 Nov 15:00 AST. The 40-epoch runs are light and may continue at low CPU afterwards.
\item \textbf{Interruptions:}
\begin{itemize}
\item Any run interrupted by the time limit, a sign-out or a restart is moved to \nolinkurl{experiments/abandoned_runs/}
with a README and reported. It is never resumed.
\item A run interrupted after its test inference started is never repeated.
\item \textbf{Clarification, 11 Nov 08:10 UTC, after the PC restart.} A run interrupted before checkpoint selection and
test inference is restarted from scratch under the same ID, as on 1--6 Nov.
\begin{itemize}
\item This applied to three 40-epoch seed-42 runs stopped in epoch 1:
\nolinkurl{experiments/abandoned_runs/2026-11-11_pc_restart_long40/}.
\item It is not a change of analysis.
\end{itemize}
\end{itemize}
\item \textbf{Hardware:} processes run hidden at Idle priority, limited to 8 logical CPUs overnight and 2 from 15:00.
Runtimes are therefore not speed measurements.
\begin{itemize}
\item \textbf{What happened (recorded 11 Nov, evening).} The six ResNet-18 runs finished overnight on 8 logical CPUs, by
11 Nov 03:00 UTC. The PC was then free until 22:00 local time, so the 40-epoch runs used 10 logical CPUs from
08:05 UTC and 11 from 08:50 UTC (about 70\% of the processor, as the user allowed). Hardware limits do not change
any result.
\end{itemize}
\item \textbf{Known limitations:}
\begin{itemize}
\item three seeds only;
\item the external-test extension is outside the frozen plan;
\item ResNet-18 sees the 2$\times$ max-pooled page (4 px/mm), not the full 8 px/mm render;
\item GPU training is not bit-reproducible, and some cuDNN backward kernels (pooling) are non-deterministic.
\end{itemize}
\end{itemize}

\subsection{Amendment of 11 November 2026 (evening): fixed-gain training with a 40-epoch budget}


Written and frozen \textbf{before} any run listed here started (11 Nov 2026, committed 23:10 local time = 19:10 UTC, commit 652b435).
Nothing below changes a completed run, or a config or hyperparameter used by an existing run. This is a
\textbf{sensitivity analysis}. It does not replace the primary results.

\textbf{Motivation.} The 40-epoch sensitivity of the morning amendment (\nolinkurl{docs/protocol_amendment_2026-11-11.md}, part A)
showed that the 20-epoch budget limited every adaptive-gain head except the ViT, by 0.012--0.019 test macro AUROC. The
paper's main finding, that fixed-gain training improves on adaptive-gain training, was measured with 20-epoch models.
\begin{itemize}
\item If fixed gain mainly speeds up convergence, the advantage could shrink when both are trained longer.
\item This amendment tests that.
\item The fixed-gain 40-epoch results have not been seen by anyone. The adaptive-gain 40-epoch results are known
(Supplementary Tables S16--S17).
\end{itemize}
\paragraph{Runs}
\begin{itemize}
\item \textbf{Heads:}
\begin{itemize}
\item CNN-only (\nolinkurl{cnn_pool});
\item CNN-LSTM (\nolinkurl{cnn_lstm});
\item CNN + temporal convolution (\nolinkurl{temporal_conv});
\item CNN + small ViT (\nolinkurl{hybrid_vit}).
These are the four architectures with fixed-gain runs; the capacity-matched head was never trained on fixed gain.
\end{itemize}
\item \textbf{Pages and seeds:} fixed-gain (2.5 mm/mV) pages; seeds \textbf{42, 43, 44}; all on the PC GPU. 12 runs in total.
\item \textbf{Configs:} \nolinkurl{experiments/configs/ptbxl_v2_long40_fg_<head>.json}. Each is a copy of the head's existing
fixed-gain config \nolinkurl{ptbxl_v2_fg_<head>.json}, with only these changes:
\begin{itemize}
\item \nolinkurl{max_epochs} 40 instead of 20, so the cosine schedule spans 40 epochs;
\item a new protocol name, result family \nolinkurl{ptbxl_v2_long40} and run IDs \nolinkurl{ptbxl_v2_long40_fg_<head>_seed<s>}.
\end{itemize}
\item \textbf{Unchanged:} everything else, including patience 4, the selection rule, threshold 0.5, one test inference per
run, the shared trainer and \nolinkurl{src/run_ptbxl_extra.py}. These are the same settings as the adaptive-gain 40-epoch
runs.
\end{itemize}
\paragraph{Analysis (fixed now)}
\begin{enumerate}
\item \textbf{Selected epoch and epochs run} per run.
\item \textbf{Primary: fixed $-$ adaptive at 40 epochs.}
\begin{itemize}
\item What is compared: macro AUROC and HYP AUROC on the PTB-XL test fold, with the 40-epoch adaptive-gain runs of
the same head and seeds.
\item Method: hierarchical seed $\times$ patient bootstrap (2,000 resamples, seed 2024), the same as the main fixed-gain
analysis.
\item Correction: Holm over the 8 comparisons (4 heads $\times$ 2 metrics).
\end{itemize}
\item \textbf{Secondary: fixed-gain 40 $-$ 20 epochs} per head (macro AUROC, seeds 42--44, Holm over 4). The 20-epoch
fixed-gain runs are all GPU-trained.
\item \textbf{Descriptive:} the 20-epoch fixed $-$ adaptive differences for the same seeds 42--44, next to the 40-epoch ones.
\item \textbf{Interpretation rule.}
\begin{itemize}
\item Macro and HYP differences stay positive for every head, with Holm-adjusted p < 0.05 where they were
significant at 20 epochs: the gain finding is robust to the training budget.
\item A difference shrinks below significance or turns negative: the manuscript reports this and qualifies the
gain finding for that head.
\item Either outcome is reported.
\end{itemize}
\item \textbf{External cohorts.} No external inference for these runs. The external tests stay with the 20-epoch models, as
stated in the paper.
\end{enumerate}
\paragraph{Execution}
\begin{itemize}
\item Four hidden lanes at Idle priority, limited to 11 logical CPUs (about 70\% of the processor), started 11 Nov
evening. The PC is free overnight.
\item Runs interrupted before checkpoint selection and test inference are moved to \nolinkurl{experiments/abandoned_runs/} with a
README and restarted from scratch under the same ID, as on 1--11 Nov. A run interrupted after its test inference
started is never repeated.
\item \textbf{Known limitations:}
\begin{itemize}
\item three seeds;
\item GPU training is not bit-reproducible;
\item seed 42 of the 20-epoch adaptive-gain runs of four heads was CPU-trained. This does not affect the primary
comparison, because all 40-epoch runs are on the GPU.
\end{itemize}
\end{itemize}

\subsection{Amendment of 15 November 2026: capacity-matched head on fixed-gain pages}


Written and frozen \textbf{before} any run listed here started (15 Nov 2026, evening, local time; the commit that adds
this file and the config precedes the first run). Nothing below changes a completed run, or a config or
hyperparameter used by an existing run. Everything produced under this amendment is labelled **Additional
analysis** in the manuscript and the Supplementary Information and is kept apart from the earlier analyses.

\textbf{Motivation.} The supervisor's review (15 Nov 2026) noted that the Supplementary Information says "the
capacity-matched head was never trained on fixed gain" and asked for it to be run "using the same five seeds and
the same evaluation settings". The request reached us as "the same five C's"; we read this as the five training
seeds 42--46, which is the only five-fold repetition in the protocol. The other possible reading, the five
diagnostic superclasses (NORM, MI, STTC, CD, HYP), is satisfied by the same runs, because they use the same five
labels. Both readings therefore lead to the runs below. The fixed-gain results of this head have not been seen by
anyone.
\paragraph{Runs}
\begin{itemize}
\item \textbf{Head:} capacity-matched CNN-only (\nolinkurl{cnn_pool_wide}, 17,065 parameters), unchanged model file
\nolinkurl{src/models/ptbxl_cnn_pool_wide.py}.
\item \textbf{Pages and seeds:} fixed-gain (2.5 mm/mV) pages, dataset \nolinkurl{ptbxl_images_fixedgain2p5_v1} (same cache and
SHA-256 binding as every other fixed-gain run); seeds \textbf{42, 43, 44, 45, 46}; all on the PC GPU. 5 runs.
\item \textbf{Config:} \nolinkurl{experiments/configs/ptbxl_v2_fg_cnn_pool_wide.json}. It is the head's adaptive-gain config
\nolinkurl{ptbxl_v2_cnn_pool_wide.json} with only the dataset fields (dataset id, cache root, cache hashes, loader,
\nolinkurl{render_gain} note) taken from the existing fixed-gain configs, a new protocol name, result family
\nolinkurl{ptbxl_v2_fixedgain} and run IDs \nolinkurl{ptbxl_v2_fg_cnn_pool_wide_seed<s>}. The \nolinkurl{training} block is identical
(checked by script): AdamW, lr 1e-3, weight decay 1e-4, batch 16, at most 20 epochs, cosine to 1e-5, patience 4,
dropout 0.2, clip 1.0.
\item \textbf{Unchanged:} selection rule (validation macro AUROC, then macro AP, then $-$BCE), threshold 0.5, one test
inference per run after the checkpoint is frozen and reloaded, the shared trainer, and the runner
\nolinkurl{src/run_ptbxl_extra.py} (the same runner as the 11 Nov amendments).
\item \textbf{Hardware note.} All five adaptive-gain runs of this head were GPU-trained, so this head's comparison has no
CPU/GPU mix.
\end{itemize}
\paragraph{Analysis (fixed now, before any result exists)}

Bootstrap: hierarchical seed $\times$ patient bootstrap on the PTB-XL test fold (seeds resampled, then patients, both
shared by the two models), seed 2024, percentile 95\% CI; seed $\times$ record bootstrap on the external cohorts. The
number of resamples and the p-value formula follow the statistical revision
\nolinkurl{docs/protocol_amendment_2026-11-15_bootstrap_inference.md}, which is frozen before any of these results is
analysed; the legacy formula (2,000 resamples) is also computed so that the new rows can be read next to the old
tables.
\begin{enumerate}
\item \textbf{Per run:} selected epoch, validation and test metrics, per-label AUROC (same table format as every run).
\item \textbf{Primary: fixed $-$ adaptive for the capacity-matched head} on the PTB-XL test fold, macro AUROC and HYP AUROC,
seeds 42--46. Holm over its own family of 2. As a sensitivity analysis it is also placed in a family of 10 with
the eight existing fixed-versus-adaptive comparisons (five heads $\times$ two metrics), and both adjusted p-values
are reported.
\item \textbf{Secondary: head pairs on fixed-gain pages} (PTB-XL test, macro AUROC, seeds 42--46): all 10 pairs among the
five heads, Holm over 10. This asks whether the sequence heads still exceed the capacity-matched head when
the gain cue is removed.
\item \textbf{External cohorts.} One inference per run and cohort (Georgia, China) through
\texttt{src.external\_test infer-added}, after the usual reproduction check of the saved validation probabilities.
Reported: external macro AUROC per run, fixed $-$ adaptive for this head (macro and HYP; Holm over 2 within each
cohort), and the 10 fixed-gain head pairs (Holm over 10). These are \textbf{outside} the frozen confirmatory
families H1 and H2 of \nolinkurl{docs/external_test_plan.md}, which are not changed.
\item \textbf{Interpretation rule.}
\begin{itemize}
\item The direction and size of fixed $-$ adaptive are reported whatever they are.
\item A difference is called an improvement only if its 95\% CI excludes zero and its Holm-adjusted p < 0.05;
otherwise it is described as a point estimate that is not statistically supported.
\item If the capacity-matched head is worse with fixed gain on either metric, the manuscript sentences that
describe fixed gain as raising the point estimates of every architecture are corrected.
\item No non-inferiority claim is made: no margin was prespecified.
\end{itemize}
\end{enumerate}
\paragraph{Execution}
\begin{itemize}
\item Hidden lanes at low priority on at most 8 of 16 logical CPUs (shared PC), started 15 Nov evening.
\item Runs interrupted before checkpoint selection and test inference are moved to \nolinkurl{experiments/abandoned_runs/} with
a README and restarted from scratch under the same ID. A run interrupted after its test inference started is
never repeated.
\item \textbf{Known limitations:} five seeds; GPU training is not bit-reproducible; 20-epoch budget (this head selected
epoch 18--20 in every adaptive-gain run, so it is probably budget-limited here as well); the request for this
analysis came after all earlier results were known, so it is not part of the prospectively frozen plan.
\end{itemize}

\subsection{Statistical revision of 15 November 2026: finite-resample bootstrap $p$-values}


Written and frozen on 15 Nov 2026 \textbf{before} any statistic was recomputed and before any result of the
fixed-gain capacity-matched runs (amendment \nolinkurl{docs/protocol_amendment_2026-11-15_fixedgain_wide.md}) existed.
It changes \textbf{how p-values are computed from the bootstrap}, for every comparison in the paper. It does not
change any model, prediction, hypothesis, comparison family, test statistic or bootstrap design.
\paragraph{What was wrong}

Every comparison uses the same routine (\nolinkurl{src/external_test.py::hierarchical_difference}): 2,000 hierarchical
bootstrap replicates $\Delta$* of the difference in mean-over-seeds AUROC, a percentile 95\% CI, and

p\_old = 2 $\cdot$ min( \#\{$\Delta$* $\leq$ 0\} / B , \#\{$\Delta$* $\geq$ 0\} / B ),   B = 2,000.
\begin{enumerate}
\item \textbf{p\_old can be exactly 0.} When no replicate crosses zero the estimate is 0, and the Holm step
(multiplication by the number of remaining hypotheses) leaves it at 0. The tables printed this as
"p < 0.001" and the Methods said "an adjusted p < 0.001 means that no bootstrap replicate crossed zero".
That statement is not supported by 2,000 resamples: a Monte Carlo p-value from B resamples cannot be shown to
be smaller than about 1/(B+1) per tail. With B = 2,000 the smallest supportable two-sided p is
2/2,001 $\approx$ 0.001, and after Holm correction over 10 comparisons it is $\approx$ 0.01 (0.008 over 8 comparisons).
All adjusted "p < 0.001" entries were therefore over-stated by up to an order of magnitude. This was raised
in the supervisor's review of 15 Nov 2026.
\item The resolution was never stated.
\end{enumerate}
\paragraph{Revised procedure (applies to every comparison)}
\begin{itemize}
\item \textbf{Resamples:} B = 10,000 (seed 2024). The first 2,000 replicates are the same random draws as before, so the
old numbers must be reproduced exactly from them; this is checked for every comparison and recorded.
\item \textbf{p-value:} the finite-resample (add-one) estimator, which can never be zero
(Davison \& Hinkley, *Bootstrap Methods and their Application*, 1997, ch. 4; Phipson \& Smyth,
Stat. Appl. Genet. Mol. Biol. 2010;9:39):
p = min\{ 1, 2 $\cdot$ min( (1 + \#\{$\Delta$* $\leq$ 0\}) / (B + 1) , (1 + \#\{$\Delta$* $\geq$ 0\}) / (B + 1) ) \}.
This is the p-value obtained by inverting the percentile interval, so "p < 0.05" and "the 95\% percentile CI
excludes zero" agree up to Monte Carlo error.
\item \textbf{Holm correction:} unchanged step-down procedure, within the same families as before (10 head pairs; 8
fixed-versus-adaptive comparisons; the families of the dated sensitivity analyses). Applied to the revised p.
\item \textbf{Resolution, always reported:} the smallest attainable raw p is 2/10,001 $\approx$ 0.0002. The smallest attainable
Holm-adjusted p in a family of m comparisons is m $\times$ 0.0002 (0.002 for m = 10; 0.0016 for m = 8; 0.001 for
m = 5; 0.0008 for m = 4; 0.0004 for m = 2). A value at this floor is printed as "$\leq$" that value and means only
that none of the 10,000 replicates crossed zero; no smaller p is claimed.
\item \textbf{Confidence intervals:} percentile 95\% intervals from the 10,000 replicates (they may differ from the old
ones in the third decimal by Monte Carlo error).
\item \textbf{Sensitivity check (reported in the supplement, not used for claims):} a null-centred bootstrap p,
p\_c = (1 + \#\{ |$\Delta$* $-$ $\Delta$̂| $\geq$ |$\Delta$̂| \}) / (B + 1), which tests H0: $\Delta$ = 0 by shifting the bootstrap distribution to
the null (Hall \& Wilson, Biometrics 1991;47:757).
\end{itemize}
\paragraph{Wording rule for the manuscript (fixed now)}
\begin{itemize}
\item A difference is called \textbf{statistically supported} only if its Holm-adjusted p < 0.05 (equivalently its
interval, after correction, is clearly away from zero). Otherwise it is reported as a \textbf{point estimate}
with its interval, and words such as "improved", "better", "outperformed" or "never worse" are not used for it.
\item No non-inferiority or equivalence claim is made anywhere: no margin was prespecified and no such test was done.
External hypothesis H2 was worded "not worse" in the frozen plan; it is evaluated as a two-sided
superiority test, as it always was computed, and is now described that way.
\end{itemize}
\paragraph{Limitations that remain (stated in the paper)}
\begin{itemize}
\item The top level of the hierarchy has only five training seeds (three in the dated sensitivity analyses). A
bootstrap over so few units cannot represent between-seed variance well, so intervals are likely somewhat too
narrow and p-values somewhat too small. The p-values are approximate.
\item The external cohorts have no patient identifiers, so records are resampled; repeated ECGs of one patient
would make those intervals too narrow.
\item The percentile-inversion p-value is an approximation, not an exact test.
\item PTB-XL test-fold comparisons remain exploratory regardless of their p-values.
\item This revision was made after the results were known. It can only weaken, not strengthen, the earlier
p-values at the floor; other p-values change by Monte Carlo error only.
\end{itemize}
\paragraph{Record keeping}
\begin{itemize}
\item The old outputs are archived byte-for-byte, with SHA-256 hashes, in
\nolinkurl{results/stats_legacy_2000resamples_20261115/} before anything is regenerated, and remain in the Git history.
\item The routine keeps the old formula as \nolinkurl{p_legacy_2000} in every output row, so each table can be read both ways.
\item The Supplementary Information gets a table of every comparison whose adjusted p changed category
(old versus revised).
\item The frozen external plan (\nolinkurl{docs/external_test_plan.md}) specified 2,000 resamples. Hypotheses, cohorts,
labels, families and the test statistic are unchanged; only the p-value estimator and the number of resamples
differ, and both versions are reported.
\end{itemize}

\subsection{Amendment of 15 November 2026: label-mapping documentation and sensitivity}


Written and frozen on 15 Nov 2026 \textbf{before} any excluded external record was rendered or given to a model and
before any sensitivity metric was computed. Known at the time of writing, from the saved inclusion tables only
(no model output): which SNOMED CT codes the excluded records carry. Nothing below changes the frozen
confirmatory analysis (\nolinkurl{docs/external_test_plan.md}), its cohorts, its predictions or its mapping.

\textbf{Motivation.} The supervisor's review (15 Nov 2026) asked for the clinical name of every SNOMED CT code, a
validation of the code-to-diagnosis mapping, and an analysis of how the mapping affects the results, because
the mapping excluded 41\% of the Georgia records (4,245 of 10,344) and 80\% of the Chinese records (36,146 of
45,152).
\paragraph{A. Clinical names and validation of the mapping}
\begin{itemize}
\item Every code in \nolinkurl{experiments/external/label_mapping.csv} gets its clinical name from the Challenge organisers'
own code tables (\nolinkurl{dx_mapping_scored.csv}, \nolinkurl{dx_mapping_unscored.csv}, repository
physionetchallenges/evaluation-2021), cross-checked against the SNOMED CT browser where reachable. The table
is saved as \nolinkurl{experiments/external/snomed_code_names.csv} with the source of every name.
\item \textbf{No cardiologist reviewed the mapping.} This is stated in the paper. The validation that can be done
without an expert is: (i) names from two independent sources agree; (ii) each assigned code's clinical meaning
is compared with the definition of its PTB-XL superclass and every mismatch is listed; (iii) the data-driven
shares are reported for every code.
\end{itemize}
\paragraph{B. Who was excluded (descriptive, from the inclusion tables)}

Per cohort: the codes carried by excluded and included records (count and share, with names), the number of
excluded records carrying only rhythm statements, label prevalence in the included cohort next to PTB-XL test,
and the positives contributed by each assigned code.
\paragraph{C. Sensitivity of the results to the mapping}

Variants (per cohort; labels eligible if $\geq$ 50 positives and $\geq$ 50 negatives, as in the frozen plan):

| Variant | Records | Mapping |
|---|---|---|
| V0 (published) | frozen cohort | 90\% rule |
| V1 strict | frozen cohort; records left without a label are dropped | share $\geq$ 0.99 |
| V2 relaxed, same records | frozen cohort | share $\geq$ 0.80 (labels added) |
| V3 all records | frozen cohort + every record excluded only for "no mapped superclass", as negatives for all labels | 90\% rule |
| V4 all records, relaxed | as V3 | share $\geq$ 0.80 |
\begin{itemize}
\item V0--V2 use the saved predictions only. V3--V4 need predictions for the unmapped records: these records are
rendered with the unchanged adaptive and fixed-gain renderers into a \textbf{separate} cache
(\nolinkurl{data/external/cache_unmapped/<cohort>/}), every completed run is reloaded and must reproduce its saved PTB-XL
validation probabilities, and each run gets exactly one inference on them
(\nolinkurl{results/external_unmapped/<cohort>/<run>/}). The frozen predictions of the included records are reused, never
recomputed.
\item Reported per variant: records, positives and prevalence per label, eligible labels, mean external macro AUROC
and per-label AUROC per model, and the H1 and H2 differences. For V3 the H1 and H2 differences get the
hierarchical bootstrap of the statistical revision (10,000 resamples, Holm within family); the other variants
are reported as point estimates.
\item \textbf{Interpretation rule.} A conclusion is called robust to the mapping if the sign of the difference is the same
in every variant and its Holm-adjusted p stays below 0.05 in V3. Any change of sign or loss of support is
reported in the main text. No causal statement about why performance differs between cohorts or variants is
made from these analyses; they show association with the cohort definition only.
\item \textbf{Known limits.} In V3/V4 unmapped records are counted as negatives although some carry statements close to a
superclass (for example T-wave abnormality, STTC share 0.88 in PTB-XL), so V3 has label noise in the
negatives; V4 reduces it for the three near-miss codes. Normal ECGs still cannot be scored as NORM, because
the organisers have no code that maps to it. The mapping remains data-driven.
\end{itemize}
\paragraph{Execution}

Rendering on at most 3 CPU workers at low priority; inference on the GPU after the fixed-gain runs of the same
day have finished. Interrupted steps are restarted; an inference that has started for a run is never repeated.
\paragraph{Addendum, 15 November 2026 (written before any sensitivity metric or any model output on the unmapped records existed)}

While collecting the clinical names (part A), the organisers' code tables showed that they score some code pairs
as the same diagnosis: right bundle branch block 59118001 = complete right bundle branch block 713427006, and
complete left bundle branch block 733534002 = left bundle branch block 164909002. The first code of each pair
never occurs in the organisers' PTB-XL encoding, so the data-driven rule could not assign it, although its
equivalent is assigned to CD. One more variant is therefore added:

| Variant | Records | Mapping |
|---|---|---|
| V5 all records, equivalents | as V3 | 90\% rule, and 59118001 and 733534002 inherit the assignment (CD) of their organiser-declared equivalents |

V5 is reported as point estimates, like V1, V2 and V4. No other code is added by judgement: for example
"left ventricular high voltage" (55827005, China only, absent from PTB-XL) stays unmapped in every variant and is
named as a limitation, because assigning it to HYP would be our own clinical decision without expert review.

\subsection{Protocol note of 15 November 2026: two descriptive additions}


Written on 15 Nov 2026 before either quantity was computed. Both add uncertainty estimates to differences whose
\textbf{point estimates are already in the paper} (Tables 1 and 4), so the direction of each is known; nothing is
trained and no inference is run. They respond to the supervisor's review (points 6 and 8), which asked that fixed
gain's lower PCA + LDA macro AUROC be acknowledged and that the signal-versus-image gap be reported separately
for the compact models and for ResNet-18. Both are post hoc and are labelled \textbf{Additional analysis}.
\begin{enumerate}
\item \textbf{PCA + LDA, fixed minus adaptive gain.} Runs \nolinkurl{ptbxl_v2_fg_lda_seed42} and \nolinkurl{ptbxl_v2_lda_wide_seed42} (the two
single fits reported in the paper). Macro AUROC and HYP AUROC on the PTB-XL test fold (patient-cluster
bootstrap) and on each external cohort (record bootstrap), 10,000 resamples, Holm over the 2 metrics within each
cohort. There is no training seed, so the interval reflects test sampling only.
\item \textbf{Image model minus signal reference.} For each image network and page type (five heads and ResNet-18, adaptive
and fixed gain: 12 groups), the difference in mean-over-seeds macro AUROC from the 1D ResNet signal reference,
on the PTB-XL test fold and each external cohort, with the hierarchical bootstrap (seeds 42--46; 42--44 for
ResNet-18) and Holm over the 12 comparisons within each cohort. Seeds of different architectures are paired only
by their number, as in the head-pair comparisons.
\end{enumerate}

Interpretation rule: a gap is described as statistically supported only if its Holm-adjusted p < 0.05; a gap whose
interval includes zero is described as "not distinguishable in this sample", not as equality (no equivalence margin
was prespecified).

\begin{thebibliography}{99}
\bibitem{collins2024tripodai} Collins, Gary S, Moons, Karel G M, Dhiman, Paula, Riley, Richard D, Beam, Andrew L, Van Calster, Ben, et al. TRIPOD+AI statement: updated guidance for reporting clinical prediction models that use regression or machine learning methods. \emph{BMJ} 385, e078378 (2024). doi:\nolinkurl{10.1136/bmj-2023-078378}
\bibitem{shivashankara2024ecgimagekit} Shivashankara, Kshama Kodthalu, Deepanshi, Mehri Shervedani, Afagh, Clifford, Gari D, Reyna, Matthew A, Sameni, Reza. ECG-Image-Kit: a synthetic image generation toolbox to facilitate deep learning-based electrocardiogram digitization. \emph{Physiological Measurement} 45, 055019 (2024). doi:\nolinkurl{10.1088/1361-6579/ad4954}
\bibitem{reyna2024challenge} Reyna, Matthew A., Deepanshi, Weigle, James Tyler, Koscova, Zuzana, Campbell, Kiersten S., Seyedi, Salman, et al. Digitization and Classification of ECG Images: The George B. Moody PhysioNet Challenge 2024. \emph{Computing in Cardiology Conference (CinC)} 51 (2024). doi:\nolinkurl{10.22489/CinC.2024.011}
\bibitem{reyna2024imagedatabase} Reyna, Matthew A., Deepanshi, Weigle, James, Koscova, Zuzana, Campbell, Kiersten, Shivashankara, Kshama Kodthalu, et al. ECG-Image-Database: A Dataset of ECG Images with Real-World Imaging and Scanning Artifacts; A Foundation for Computerized ECG Image Digitization and Analysis. \emph{arXiv preprint; arXiv:2409.16612} (2024). doi:\nolinkurl{10.48550/arXiv.2409.16612}
\bibitem{sangha2022images} Sangha, Veer, Mortazavi, Bobak J., Haimovich, Adrian D., Ribeiro, Ant\^onio H., Brandt, Cynthia A., Jacoby, Daniel L., et al. Automated multilabel diagnosis on electrocardiographic images and signals. \emph{Nature Communications} 13, 1583 (2022). doi:\nolinkurl{10.1038/s41467-022-29153-3}
\bibitem{khan2021ecgimages} Khan, Ali Haider, Hussain, Muzammil, Malik, Muhammad Kamran. ECG Images dataset of Cardiac and COVID-19 Patients. \emph{Data in Brief} 34, 106762 (2021). doi:\nolinkurl{10.1016/j.dib.2021.106762}
\end{thebibliography}
\end{document}
